\subsection{INTEGRATION FORMULA OF H. WEYL}

Let e be the identity element of G and \(\bar{e}\) its class in G/T. Identify the tangent space of G at e with g, the tangent space of T at e with t and the tangent space of G/T at \(\bar{e}\) with g/t. Denote by \((u,v)\mapsto u\cap v\) the R-bilinear map

\[
\operatorname{Alt} ^ {n - r} (\mathfrak {g} / \mathfrak {t}) \times \operatorname{Alt} ^ {r} (\mathfrak {t}) \rightarrow \operatorname{Alt} ^ {n} (\mathfrak {g})
\]

defined in number 1.

Recall (Chap. III, §3, no. 13, Prop. 50) that the map \(\omega \mapsto \omega(e)\) is an isomorphism from the vector space of left-invariant differential forms of degree \(n\) (resp. \(r\) ) on G (resp. T) to the space \(\mathrm{Alt}^n(\mathfrak{g})\) (resp. \(\mathrm{Alt}^r(t)\) ). Further, observe that, since every connected compact subgroup of \(\mathbf{R}^*\) reduces to the identity element, \(\det \mathrm{Ad}g = 1\) for all \(g \in \mathrm{G}\) , so that the left-invariant differential forms of degree \(n\) on G are also right-invariant and invariant under inner automorphisms (Chap. III, §3, no. 16, Cor. of Prop. 54): we shall speak simply of G-invariant differential forms from now on.

Similarly, it follows from Chap. III, §3, no. 16, Prop. 56 and the preceding that the map \(\omega \mapsto \omega (\bar{e})\) is an isomorphism from the space of G-invariant differential forms of degree \(n - r\) on \(\mathrm{G / T}\) to the space \(\mathrm{Alt}^{n - r}(\mathfrak{g} / \mathfrak{t})\) .

If \(\omega_{G/T}\) is a G-invariant differential form of degree n - r on G/T, and \(\omega_{T}\) an invariant differential form of degree r on T, denote by \(\omega_{G/T} \cap \omega_{T}\) the unique invariant differential form of degree n on G such that

\[
(\omega_ {\mathrm{G/T}} \cap \omega_ {\mathrm{T}}) (e) = \omega_ {\mathrm{G/T}} (\bar {e}) \cap \omega_ {\mathrm{T}} (e).
\]

Recall finally that \(f : (\mathrm{G}/\mathrm{T}) \times \mathrm{T} \to \mathrm{G}\) denotes the morphism of manifolds induced by the map \((g, t) \mapsto gtg^{-1}\) from \(G \times T\) to G by passage to the quotient (§5, no. 4). If \(\alpha\) and \(\beta\) are differential forms on G/T and T, respectively, denote simply by \(\alpha \wedge \beta\) the form \(pr_{1}^{*}\alpha \wedge pr_{2}^{*}\beta\) on \((\mathrm{G}/\mathrm{T}) \times \mathrm{T}\) .

For \(t \in \mathrm{T}\) , denote by \(\operatorname{Ad}_{\mathfrak{g} / \mathfrak{t}}(t)\) the endomorphism of \(\mathfrak{g} / \mathfrak{t}\) induced by \(\operatorname{Ad} t\) by passage to the quotient. Put

\[
\delta_ {\mathrm{G}} (t) = \det (\mathrm{Ad} _ {\mathfrak {g} / \mathfrak {t}} (t) - 1) = \prod_ {\alpha \in \mathrm{R} (\mathrm{G}, \mathrm{T})} (t ^ {\alpha} - 1).\tag{3}
\]

Let \(x \in \mathfrak{t}\) and \(\alpha \in \mathrm{R}(\mathrm{G},\mathrm{T})\) ; denote by \(\hat{\alpha}\) the element \((2\pi i)^{-1}\delta (\alpha)\) of \(\mathfrak{t}^*\) , so that

\[
((\exp x) ^ {\alpha} - 1) ((\exp x) ^ {- \alpha} - 1) = (e ^ {2 \pi i \hat {\alpha} (x)} - 1) (e ^ {- 2 \pi i \hat {\alpha} (x)} - 1) = 4 \sin^ {2} \pi \hat {\alpha} (x).
\]

If \(\mathrm{R}_{+}(\mathrm{G},\mathrm{T})\) denotes the set of positive roots of \(\mathrm{R}(\mathrm{G},\mathrm{T})\) relative to a basis B, we have

\[
\delta_ {\mathrm{G}} (\exp x) = \prod_ {\alpha \in \mathrm{R} _ {+} (\mathrm{G}, \mathrm{T})} 4 \sin^ {2} \pi \hat {\alpha} (x),
\]

so, in particular, \(\delta_{\mathrm{G}}(t)>0\) for all \(t\in T_{r}\) . We remark also that \(\delta_{\mathrm{G}}(t)=\delta_{\mathrm{G}}(t^{-1})\) for \(t\in T\) .

PROPOSITION 1. Let \(\omega_{\mathrm{G}}, \omega_{\mathrm{G/T}}\) and \(\omega_{\mathrm{T}}\) be invariant differential forms on G, G/T and T, respectively, of respective degrees \(n, n - r\) and \(r\) . If \(\omega_{\mathrm{G}} = \omega_{\mathrm{G/T}} \cap \omega_{\mathrm{T}}\) , then

\[
f ^ {*} (\omega_ {\mathrm{G}}) = \omega_ {\mathrm{G/T}} \wedge \delta_ {\mathrm{G}} \omega_ {\mathrm{T}}.
\]

Clearly we can assume that \(\omega_{G/T}\) and \(\omega_{T}\) are non-zero; then the differential form \((u,t)\mapsto\omega_{\mathrm{G/T}}(u)\wedge\omega_{\mathrm{T}}(t)\) on \((\mathrm{G/T})\times\mathrm{T}\) is of degree n and everywhere non-zero; hence there exists a numerical function \(\delta\) on \((\mathrm{G/T})\times\mathrm{T}\) such that

\[
f ^ {*} (\omega_ {\mathrm{G}}) (u, t) = \delta (u, t) \omega_ {\mathrm{G/T}} (u) \wedge \omega_ {\mathrm{T}} (t).
\]

Observe now that, for \(h \in \mathrm{G}\) , \(u \in \mathrm{G} / \mathrm{T}\) , \(t \in \mathrm{T}\) , we have \(f(h.u,t) = (\operatorname{Int} h)f(u,t)\) ; since \(\omega_{\mathrm{G}}\) is invariant under inner automorphisms, it follows immediately that \(\delta(h.u,t) = \delta(u,t)\) , so \(\delta(u,t) = \delta(\bar{e},t)\) .

Denote by \(p:\mathfrak{g}\to \mathfrak{g} / \mathfrak{t}\) the quotient map and by \(\varphi :\mathfrak{g} / \mathfrak{t}\rightarrow \mathfrak{g}\) the map defined by

\[
\varphi (p (X)) = (\operatorname{Ad} t ^ {- 1}) X - X \text { for } X \in \mathfrak {g};
\]

recall (§5, no. 4, Lemma 4) that the tangent map

\[
\mathrm{T} _ {(e, t)} (f): \mathrm{T} _ {e} (\mathrm{G} / \mathrm{T}) \times \mathrm{T} _ {t} (\mathrm{T}) \to \mathrm{T} _ {t} (\mathrm{G})
\]

takes \((z,tH)\) to \(t(\varphi (z) + H)\) for \(z\in \mathfrak{g} / \mathfrak{t},H\in \mathfrak{t}\) .

Let \(z_{1},\ldots,z_{n-r}\) be elements of g/t, \(H_{1},\ldots,H_{r}\) elements of t. Then

\[
\begin{array}{l} f ^ {*} \omega_ {\mathrm{G}} (\bar {e}, t) (z _ {1}, \ldots , z _ {n - r}, t H _ {1}, \ldots , t H _ {r}) \\ \quad = \omega_ {\mathrm{G}} (t) (t \varphi (z _ {1}), \ldots , t \varphi (z _ {n - r}), t H _ {1}, \ldots , t H _ {r}) \text {by the calculation of T} _ {(\bar {e}, t)} (f) \\ \quad = \omega_ {\mathrm{G}} (e) (\varphi (z _ {1}), \ldots , \varphi (z _ {n - r}), H _ {1}, \ldots , H _ {r}) \text {since} \omega_ {\mathrm{G}} \text {is invariant} \\ \quad = \omega_ {\mathrm{G/T}} (\bar {e}) (p \varphi (z _ {1}), \ldots , p \varphi (z _ {n - r})) \cdot \omega_ {\mathrm{T}} (e) (H _ {1}, \ldots , H _ {r}) \text {(no. 1, Lemma 1)} \\ \quad = \det (p \varphi) \omega_ {\mathrm{G/T}} (\bar {e}) (z _ {1}, \ldots , z _ {n - r}). \omega_ {\mathrm{T}} (e) (H _ {1}, \ldots , H _ {r}) \\ \quad = \delta_ {\mathrm{G}} (t) \omega_ {\mathrm{G/T}} (\bar {e}) (z _ {1}, \ldots , z _ {n - r}). \omega_ {\mathrm{T}} (t) (t H _ {1}, \ldots , t H _ {r}) \\ \qquad \qquad \qquad \qquad \qquad \qquad \text {since} \omega_ {\mathrm{T}} \text {is invariant} \\ \quad = \delta_ {\mathrm{G}} (t) (\omega_ {\mathrm{G/T}} \wedge \omega_ {\mathrm{T}}) (\bar {e}, t) (z _ {1}, \ldots , z _ {n - r}, t H _ {1}, \ldots , t H _ {r}), \end{array}
\]

so \(f^{*}\omega_{\mathrm{G}}(\bar{e},t) = \delta_{\mathrm{G}}(t)(\omega_{\mathrm{G / T}}\wedge \omega_{\mathrm{T}})(\bar{e},t)\) ; thus, \(\delta (\bar {e},t) = \delta_{\mathrm{G}}(t)\) , hence the proposition.

Give the manifolds G, T and G/T the orientations defined by the forms \(\omega_{G}, \omega_{T}\) and \(\omega_{G/T}\) , respectively. These forms define invariant measures on G, T and G/T (Chap. III, §3, no. 16, Props. 55 and 56), also denoted by \(\omega_{G}, \omega_{T}\) and \(\omega_{G/T}\) .

Lemma 2. If \(\omega_{\mathrm{G}} = \omega_{\mathrm{G / T}}\cap \omega_{\mathrm{T}}\) , then

\[
\int_ {\mathrm{G}} \omega_ {\mathrm{G}} = \int_ {\mathrm{G/T}} \omega_ {\mathrm{G/T}}. \int_ {\mathrm{T}} \omega_ {\mathrm{T}}.
\]

Denote by \(\pi\) the canonical morphism from G to G/T. Let \(g \in G\) , and let \(t_{1}, \ldots, t_{n-r}\) be elements of \(\mathrm{T}_{\pi(g)}(\mathrm{G}/\mathrm{T})\) . Identify the fibre \(\pi^{-1}(\pi(g)) = g\mathrm{T}\) with T by the translation \(\gamma(g)\) . The relation \(\omega_{G} = \omega_{G/T} \cap \omega_{T}\) now implies the equality (Differentiable and Analytic Manifolds, Results, 11.4.5):

\[
\omega_ {\mathrm{G} \sqcup} (t _ {1}, \dots , t _ {n - r}) = (\omega_ {\mathrm{G} / \mathrm{T}} (t _ {1}, \dots , t _ {n - r})) \omega_ {\mathrm{T}}.
\]

Thus \(\int_{\pi}\omega_{\mathrm{G}} = \left(\int_{\mathrm{T}}\omega_{\mathrm{T}}\right)\omega_{\mathrm{G / T}}\) (Differentiable and Analytic Manifolds, Results, 11.4.6), and

\[
\int_ {\mathrm{G}} \omega_ {\mathrm{G}} = \int_ {\mathrm{G/T}} \int_ {\pi} \omega_ {\mathrm{G}} = \int_ {\mathrm{T}} \omega_ {\mathrm{T}}. \int_ {\mathrm{G/T}} \omega_ {\mathrm{G/T}}
\]

(Differentiable and Analytic Manifolds, Results, 11.4.8).

Lemma 3. The inverse image on \((\mathrm{G}/\mathrm{T}) \times \mathrm{T}_{r}\) of the measure dg on \(G_{r}\) under the local homeomorphism \(f_{r}\) (Integration, Chap. V, §6, no. 6) is the measure \(\mu \otimes \delta_{G} dt\) , where \(\mu\) is the unique G-invariant measure on G/T of total mass 1.

Choose an invariant differential form \(\omega_{T}\) (resp. \(\omega_{G/T}\) ) on T (resp. G/T) of maximal degree, such that the measure defined by \(\omega_{T}\) (resp. \(\omega_{G/T}\) ) is equal to dt (resp. \(\mu\) ). Put \(\omega_{G} = \omega_{G/T} \cap \omega_{T}\) . Lemma 2 implies that the measure defined by \(\omega_{G}\) is equal to dg. Let U be an open subset of (G/T) × T \(_{r}\) such that \(f_{r}\) induces an isomorphism from U to an open subset V of G \(_{r}\) . Let \(\varphi\) be a continuous function with compact support in V; denote also by \(\varphi\) the extension of \(\varphi\) to G \(_{r}\) which vanishes outside V. We have

\[
\begin{array}{r l} & {\int_ {\mathrm{V}} \varphi d g = \int_ {\mathrm{V}} \varphi \omega_ {\mathrm{G}} = \int_ {\mathrm{U}} (\varphi \circ f _ {r}) f _ {r} ^ {*} (\omega_ {\mathrm{G}})} \\ & {\qquad = \int_ {\mathrm{U}} (\varphi \circ f _ {r}) \omega_ {\mathrm{G/T}} \wedge \delta_ {\mathrm{G}} \omega_ {\mathrm{T}} \quad (\text {Prop. 1})} \\ & {\qquad = \int_ {\mathrm{U}} (\varphi \circ f _ {r}) d \mu . \delta_ {\mathrm{G}} d t,} \end{array}
\]

hence the lemma.

THEOREM 1 (H. Weyl). The measure \(dg\) on \(G\) is the image under the map \((g, t) \mapsto gtg^{-1}\) from \(G \times T\) to \(G\) of the measure \(dg \otimes \frac{1}{w(G)}\delta_G dt\) , where

\[
\delta_ {\mathrm{G}} (t) = \det (\operatorname{Ad} _ {\mathfrak {g} / \mathfrak {t}} (t) - 1) = \prod_ {\alpha \in \mathrm{R} (\mathrm{G}, \mathrm{T})} (t ^ {\alpha} - 1).
\]

Equivalently (Integration, Chap. V, §6, no. 3, Prop. 4), \(dg\) is the image under the map \(f: (\mathrm{G} / \mathrm{T}) \times \mathrm{T} \to \mathrm{G}\) of the measure \(\mu \otimes \frac{1}{w(\mathrm{G})}\delta_{\mathrm{G}}dt\) .

We prove the last assertion. It follows from §5, no. 1 and Differentiable and Analytic Manifolds, Results, 10.1.3 c) that G-G \(_{r}\) is negligible in G and T-T \(_{r}\) is negligible in T. Further, the map \(f_{r}\) makes (G/T) × T \(_{r}\) a principal covering of G \(_{r}\) , with group W (§5, no. 4, Prop. 4 b)). The theorem now follows from Lemma 3 and Integration, Chap. V, §6, no. 6, Prop. 11.

COROLLARY 1. (i) Let \(\varphi\) be an integrable function on G with values in a Banach space or in \(\overline{\mathbf{R}}\) . For almost all \(t \in \mathrm{T}\) , the function \(g \mapsto \varphi(gtg^{-1})\) on G is integrable for \(dg\) . The function \(t \mapsto \delta_{\mathrm{G}}(t) \int_{\mathrm{G}} \varphi(gtg^{-1}) dg\) is integrable on T, and we have

\[
\int_ {\mathrm{G}} \varphi (g) d g = \frac {1}{w (\mathrm{G})} \int_ {\mathrm{T}} \left(\int_ {\mathrm{G}} \varphi (g t g ^ {- 1}) d g\right) \delta_ {\mathrm{G}} (t) d t\tag{4}
\]

(``Hermann Weyl's integration formula'').

\begin{enumerate}
\def\labelenumi{(\roman{enumi})}
\setcounter{enumi}{1}
\tightlist
\item
  Let \(\varphi\) be a positive measurable function on \(\mathbf{G}\) . For almost all \(t \in \mathrm{T}\) , the function \(g \mapsto \varphi(gtg^{-1})\) on \(\mathbf{G}\) is measurable. The function \(t \mapsto \int_{\mathbf{G}}^{*} \varphi(gtg^{-1}) dg\) on \(\mathrm{T}\) is measurable, and we have
\end{enumerate}

\[
\int_ {\mathrm{G}} ^ {*} \varphi (g) d g = \frac {1}{w (\mathrm{G})} \int_ {\mathrm{T}} ^ {*} \left(\int_ {\mathrm{G}} ^ {*} \varphi (g t g ^ {- 1}) d g\right) \delta_ {\mathrm{G}} (t) d t.\tag{5}
\]

Since the map f is induced by passage to the quotient from the map \((g, t) \mapsto gtg^{-1}\) from \(G \times T \to G\) , it suffices to apply Integration, Chap. V, §5, 6, 8 and Integration, Chap. VII, §2.

COROLLARY 2. Let \(\varphi\) be a central function on \(\mathbf{G}\) (that is, such that \(\varphi(gh) = \varphi(hg)\) for all \(g\) and \(h\) in \(\mathbf{G}\) ) with values in a Banach space or in \(\overline{\mathbf{R}}\) .

\begin{enumerate}
\def\labelenumi{\alph{enumi})}
\item
  \(\varphi\) is measurable if and only if its restriction to T is measurable.
\item
  \(\varphi\) is integrable if and only if the function \((\varphi |\mathrm{T})\delta_{\mathrm{G}}\) is integrable on \(\mathrm{T}\) , and in that case we have
\end{enumerate}

\[
\int_ {\mathrm{G}} \varphi (g) d g = \frac {1}{w (\mathrm{G})} \int_ {\mathrm{T}} \varphi (t) \delta_ {\mathrm{G}} (t) d t.\tag{6}
\]

Denote by \(p: \mathrm{G} / \mathrm{T} \times \mathrm{T} \to \mathrm{T}\) the second projection. We have \(\varphi \circ f = (\varphi | \mathrm{T}) \circ p\) ; further, the image under \(p\) of the measure \(\mu \otimes \frac{1}{w(\mathrm{G})} \delta_{\mathrm{G}} dt\) is \(\frac{1}{w(\mathrm{G})} \delta_{\mathrm{G}} dt\) . The corollary now follows from Th. 1 above and Th. 1 of Integration, Chap. V, §6, no. 2, applied to the two proper maps \(f\) and \(p\) .

COROLLARY 3. Let H be a connected closed subgroup of G containing T, h its Lie algebra, and dh the Haar measure on H of total mass 1. Let \(\varphi\) be an integrable central function on G, with values in a Banach space or in R. Then the function \(h \mapsto \varphi(h) \det(\mathrm{Ad}_{\mathfrak{g}/\mathfrak{h}}(h) - 1)\) is integrable and central on H and we have

\[
\int_ {\mathrm{G}} \varphi (g) d g = \frac {w (\mathrm{H})}{w (\mathrm{G})} \int_ {\mathrm{H}} \varphi (h) \mathrm{det} (\mathrm{Ad} _ {\mathfrak {g} / \mathfrak {h}} (h) - 1) d h.\tag{7}
\]

Indeed, the function \(h \mapsto \varphi(h) \det(\mathrm{Ad}_{\mathfrak{g}/\mathfrak{h}}(h) - 1)\) is a central function on H whose restriction to T is the function \(t \mapsto \varphi(t) \delta_{\mathrm{G}}(t) \delta_{\mathrm{H}}(t)^{-1}\) . Thus, the corollary follows from Cor. 2 applied to G and to H.

Remarks. 1) If we take \(\varphi = 1\) in Cor. 3, we obtain

\[
\int_ {\mathrm{H}} \det (\mathrm{Ad} _ {\mathfrak {g} / \mathfrak {h}} (h) - 1) d h = w (\mathrm{G}) / w (\mathrm{H})\tag{8}
\]

and in particular

\[
\int_ {\mathrm{T}} \delta_ {\mathrm{G}} (t) d t = w (\mathrm{G}).\tag{9}
\]

\begin{enumerate}
\def\labelenumi{\arabic{enumi})}
\setcounter{enumi}{1}
\tightlist
\item
  Let \(\nu\) be the measure on the quotient T/W defined by
\end{enumerate}

\[
\int_ {\mathrm{T/W}} \psi (\tau) d \nu (\tau) = \frac {1}{w (\mathrm{G})} \int_ {\mathrm{T}} \psi (\pi (t)) \delta_ {\mathrm{G}} (t) d t,
\]

where \(\pi\) denotes the canonical projection of T onto T/W. Cor. 2 means that the homeomorphism \(T / W \to G / \operatorname{Int}(G)\) (§2, no. 5, Cor. 1 of Prop. 5)

transports the measure \(\nu\) to the image of the measure dg under the canonical projection \(\mathrm{G} \rightarrow \mathrm{G}/\mathrm{Int}(\mathrm{G})\) .

\begin{enumerate}
\def\labelenumi{\arabic{enumi})}
\setcounter{enumi}{2}
\tightlist
\item
  Assume that G is simply-connected. Let A be an alcove of t, and dx the Haar measure on t such that \(\int_{A} dx = 1\) . Then the measure \(\nu\) can also be obtained by transporting the measure \(\frac{1}{w(G)} \prod_{\alpha \in R_{+}(G,T)} 4 \sin^{2} \pi \hat{\alpha}(x) dx\) on \(\overline{A}\) by the homeomorphism \(\overline{A} \rightarrow T/W\) (§5, no. 2, Cor. 1 of Prop. 2).
\end{enumerate}

Example. Take G to be the group \(\mathbf{SU}(2,\mathbf{C})\) and T to be the subgroup of diagonal matrices ( \(\S3\) , no. 6); identify t with R by the choice of basis \(\{iH\}\) of t (loc. cit.). Put A = 10, \(\pi\) ; this is an alcove of t. The interval \(\overline{A} = [0, \pi]\) can be identified with the space of conjugacy classes of G, the element \(\theta\) of \(\overline{A}\) corresponding to the conjugacy class of \(\begin{pmatrix} e^{i\theta} & 0 \\ 0 & e^{-i\theta} \end{pmatrix}\) . Let \(d\theta\) be Lebesgue measure on \([0, \pi]\) ; it follows from the preceding that the image on \(\overline{A}\) of the Haar measure on G is the measure \(\frac{2}{\pi}\sin^{2}\theta d\theta\) .

\subsection{INTEGRATION ON LIE ALGEBRAS}

PROPOSITION 2. Let H be a (real) Lie group of dimension m, h its Lie algebra. Let \(\omega_{H}\) be a right-invariant differential form of degree m on H, and let \(\omega_{h}\) be the translation-invariant differential form on h, of degree m, that coincides with \(\omega_{\mathrm{H}}(e)\) at the origin. We have

\[
(\exp_ {\mathrm{H}}) ^ {*} \omega_ {\mathrm{H}} = \lambda_ {\mathfrak {h}} \omega_ {\mathfrak {h}}\tag{10}
\]

where \(\lambda_{\mathfrak{h}}\) is the Ad(H)-invariant function on \(\mathfrak{h}\) such that

\[
\lambda_ {\mathfrak {h}} (x) = \det \left(\sum_ {p \geq 0} \frac {1}{(p + 1) !} (\operatorname{ad} x) ^ {p}\right) \quad \text { for } x \in \mathfrak {h}.
\]

Let \(x, x_1, \ldots, x_m\) be elements of \(\mathfrak{h}\) . We have

\[
(\exp^ {*} \omega_ {\mathrm{H}}) _ {x} (x _ {1}, \dots , x _ {m}) = (\omega_ {\mathrm{H}} (\exp x)) (\mathrm{T} _ {x} (\exp) (x _ {1}), \dots , \mathrm{T} _ {x} (\exp) (x _ {m})).
\]

Denote by \(\varpi(x) : \mathfrak{h} \to \mathfrak{h}\) the right differential of the exponential at \(x\) (Chap. III, §3, no. 17, Def. 8); by definition,

\[
\mathrm{T} _ {x} (\exp) (y). (\exp x) ^ {- 1} = \varpi (x). y \quad \text { for   all } y \in \mathfrak {h}.
\]

The form \(\omega_{H}\) being right invariant, we obtain

\[
\begin{array}{r l} & (\omega_ {\mathrm{H}} (\exp x)) (\mathrm{T} _ {x} (\exp) (x _ {1}), \ldots , \mathrm{T} _ {x} (\exp) (x _ {m})) \\ & \quad = \omega_ {\mathrm{H}} (e) (\varpi (x). x _ {1}, \ldots , \varpi (x). x _ {m}) = (\det \varpi (x)) \omega_ {\mathfrak {h}} (x _ {1}, \ldots , x _ {m}); \end{array}
\]

thus, \(\exp^{*}\omega_{\mathrm{H}} = \lambda_{\mathfrak{h}}\omega_{\mathfrak{h}}\) , with \(\lambda_{\mathfrak{h}}(x) = \det \varpi(x) = \det \frac{\exp\operatorname{ad} x - 1}{\operatorname{ad} x}\) (Chap. III, §6, no. 4, Prop. 12).

Let \(h \in H\) ; since Ad h is an automorphism of h, we have

\[
\operatorname{ad} \left((\operatorname{Ad} h) (x)\right) = \operatorname{Ad} h \circ \operatorname{Ad} x \circ (\operatorname{Ad} h) ^ {- 1},
\]

so \(\lambda_{\mathfrak{h}}((\mathrm{Ad}h)(x)) = \lambda_{\mathfrak{h}}(x)\) . Thus, the function \(\lambda_{\mathfrak{h}}\) is invariant under \(\mathrm{Ad(H)}\) ; this completes the proof of the proposition.

Remark. Consider the function \(\lambda_{g}\) associated to a compact Lie group G; in view of §2, no. 1, Th. 1, to calculate \(\lambda_{g}\) it suffices to know its values on t. But, for \(x \in t\) , the endomorphism ad x of g is semi-simple, and has eigenvalues 0 (with multiplicity r) and, for all \(\alpha \in \mathrm{R}(\mathrm{G}, \mathrm{T})\) , \(\delta(\alpha)(x)\) (with multiplicity 1). It follows immediately that

\[
\lambda_ {\mathfrak {g}} (x) = \prod_ {\alpha \in \mathrm{R} (\mathrm{G}, \mathrm{T})} \frac {e ^ {\delta (\alpha) (x)} - 1}{\delta (\alpha) (x)} = \frac {\delta_ {\mathfrak {g}} (x)}{\pi_ {\mathfrak {g}} (x)}\tag{11}
\]

\[
\text { with } \delta_ {\mathfrak {g}} (x) = \delta_ {\mathrm{G}} (\exp x) \text { and } \pi_ {\mathfrak {g}} (x) = \prod_ {\alpha \in \mathrm{R} (\mathrm{G}, \mathrm{T})} \delta (\alpha) (x) = \det \operatorname{ad} _ {\mathfrak {g} / \mathfrak {t}} (x).
\]

Let \(\omega_{\mathrm{G / T}}\) be an invariant differential form of degree \(n - r\) on \(\mathrm{G / T}\) and \(\omega_{\mathfrak{t}}\) a translation-invariant differential form of degree \(r\) on \(\mathfrak{t}\) . With the notation of no. 1, denote by \(\omega_{\mathrm{G / T}} \cap \omega_{\mathfrak{t}}\) the unique translation-invariant differential form \(\omega_{\mathfrak{g}}\) of degree \(n\) on \(\mathfrak{g}\) such that \(\omega_{\mathfrak{g}}(0) = \omega_{\mathrm{G / T}}(\bar{e}) \cap \omega_{\mathfrak{t}}(0)\) .

Finally, denote by \(\psi:(\mathrm{G}/\mathrm{T})\times\mathfrak{t}\to\mathfrak{g}\) the morphism of manifolds induced by the map \((g,x)\mapsto(\operatorname{Ad}g)(x)\) from \(G\times t\) to g by passage to the quotient.

PROPOSITION 3. Let \(\omega_{\mathfrak{g}}\) , \(\omega_{\mathfrak{t}}\) , \(\omega_{\mathrm{G / T}}\) be invariant differential forms on \(\mathfrak{g}\) , \(\mathfrak{t}\) , \(\mathrm{G / T}\) , respectively, of respective degrees \(n, r, n - r\) . If \(\omega_{\mathfrak{g}} = \omega_{\mathrm{G / T}} \cap \omega_{\mathfrak{t}}\) , we have

\[
\psi^ {*} \omega_ {\mathfrak {g}} = \omega_ {\mathrm{G/T}} \wedge \pi_ {\mathfrak {g}} \omega_ {\mathfrak {t}}
\]

where \(\pi_{\mathfrak{g}}\) is the function on \(\mathfrak{t}\) defined by \(\pi_{\mathfrak{g}}(x) = \prod_{\alpha \in \mathrm{R}(\mathrm{G},\mathrm{T})}\delta (\alpha)(x)\) .

Denote by \(\omega_{G}\) (resp. \(\omega_{T}\) ) the invariant differential form of maximum degree on G (resp. T) that coincides with \(\omega_{g}\) (resp. \(\omega_{t}\) ) at the origin. Consider the commutative diagram

\[
\begin{array}{c c c} (\mathrm{G} / \mathrm{T}) \times \mathfrak {t} & \stackrel {{\psi}} {{\longrightarrow}} & \mathfrak {g} \\ \Big \downarrow^ {(\mathrm{Id}, \exp_ {\mathrm{T}})} & & \Big \downarrow^ {\exp_ {\mathrm{G}}} . \\ (\mathrm{G} / \mathrm{T}) \times \mathrm{T} & \stackrel {{f}} {{\longrightarrow}} & \mathrm{G} \end{array}
\]

In view of Prop. 1 of no. 2 and the relation \(\exp_{T}^{*}\omega_{T} = \omega_{t}\) , we deduce the equality

\[
\psi^ {*} \mathrm{exp} _ {\mathrm{G}} ^ {*} \omega_ {\mathrm{G}} = \omega_ {\mathrm{G/T}} \wedge \delta_ {\mathfrak {g}} \omega_ {\mathfrak {t}}.
\]

By Prop. 2, \(\psi^{*}\exp_{\mathrm{G}}^{*}\omega_{\mathrm{G}} = (\psi^{*}\lambda_{\mathfrak{g}})\psi^{*}\omega_{\mathfrak{g}}\) . Since the function \(\lambda_{\mathfrak{g}}\) is invariant under \(\operatorname{Ad}(\mathrm{G})\) , we have

\[
(\psi^ {*} \lambda_ {\mathfrak {g}}) (\bar {g}, x) = \lambda_ {\mathfrak {g}} (x) = \frac {\delta_ {\mathfrak {g}} (x)}{\pi_ {\mathfrak {g}} (x)} \quad \text { for } \bar {g} \in \mathrm{G/T}, x \in \mathfrak {t}.
\]

It follows that the forms \(\psi^{*}\omega_{\mathrm{G}}(\bar{g},x)\) and \(\omega_{\mathrm{G / T}}(\bar{g})\wedge \pi_{\mathfrak{g}}(x)\omega_{\mathfrak{t}}(x)\) coincide where \(\delta_{\mathfrak{g}}(x)\) is non-zero, that is on the dense open subset \((\mathrm{G / T})\times \mathfrak{t}_r\) ; thus, they are equal, hence the proposition.

Choose invariant differential forms \(\omega_{G}\) on G and \(\omega_{T}\) on T, of maximum degree, such that \(|\omega_{G}| = dg\) and \(|\omega_{T}| = dt\) ; denote by \(\omega_{g}\) (resp. \(\omega_{t}\) ) the translation-invariant differential form on g (resp. t) that coincides with \(\omega_{\mathrm{G}}(e)\) (resp. \(\omega_{\mathrm{T}}(e)\) ) at the origin, and dz (resp. dx) the Haar measure \(|\omega_{g}|\) (resp. \(|\omega_{t}|\) ). Reasoning as in no. 2, mutatis mutandis, gives the following proposition:

PROPOSITION 4. The measure \(dz\) on \(\mathfrak{g}\) is the image under the proper map \((g,x)\mapsto (\operatorname {Ad}g)(x)\) from \(\mathrm{G}\times \mathfrak{t}\) to \(\mathfrak{g}\) of the measure \(dg\otimes \frac{1}{w(\mathrm{G})}\pi_{\mathfrak{g}}dx\) .

We leave to the reader the statement and proof of the analogues of Cor. 1 to 3 and Remarks 1 to 3 of no. 2. For example, let \(\varphi\) be an integrable function on \(\mathfrak{g}\) (with values in a Banach space or \(\overline{\mathbf{R}}\) ); then

\[
\int_ {\mathfrak {g}} \varphi (z) d z = \frac {1}{w (\mathrm{G})} \int_ {\mathfrak {t}} \left(\int_ {\mathrm{G}} \varphi ((\mathrm{Ad} g) x) d g\right) \pi_ {\mathfrak {g}} (x) d x,\tag{12}
\]

and, in particular, if \(\varphi\) is invariant under \(\operatorname{Ad}(\mathbf{G})\) ,

\[
\int_ {\mathfrak {g}} \varphi (z) d z = \frac {1}{w (\mathrm{G})} \int_ {\mathfrak {t}} \varphi (x) \pi_ {\mathfrak {g}} (x) d x.\tag{13}
\]

\subsection{INTEGRATION OF SECTIONS OF A VECTOR BUNDLE}

In this number and the next, we denote by X a real manifold of class \(C^{r}\) \((1 \leq r \leq \infty)\) , locally of finite dimension.

Let Y be a manifold of class \(C^{r}\) . If \(r < \infty\) , consider the map \(f \mapsto j^{r}(f)\) from \(\mathcal{C}^{r}(\mathrm{X};\mathrm{Y})\) to \(\mathcal{C}(\mathrm{X};\mathrm{J}^{r}(\mathrm{X},\mathrm{Y}))\) (Differentiable and Analytic Manifolds, Results, 12.3.7). The inverse image under this map of the topology of compact convergence on \(\mathcal{C}(\mathrm{X};\mathrm{J}^{r}(\mathrm{X},\mathrm{Y}))\) is called the topology of compact \(C^{r}\) -convergence on \(\mathcal{C}^{r}(\mathrm{X};\mathrm{Y})\) ; it is the upper bound of the topologies of uniform \(C^{r}\) -convergence on K (Differentiable and Analytic Manifolds, Results, 12.3.10), where K runs through the set of compact subsets of X.

When \(r = \infty\) , we call the topology of compact \(C^{\infty}\) -convergence on \(\mathcal{C}^{\infty}(X;Y)\) the upper bound of the topologies of compact \(C^{k}\) -convergence, in other words the coarsest topology for which the canonical injections \(\mathcal{C}^{\infty}(X;Y) \to \mathcal{C}^{k}(X;Y)\) are continuous for \(0 \leq k < \infty\) .

Let \(\mathrm{E}\) be a real vector bundle with base \(\mathrm{X}\) , of class \(\mathrm{C}^r\) , and let \(\mathcal{S}^r(\mathrm{X};\mathrm{E})\) be the vector space of sections of \(\mathrm{E}\) of class \(\mathrm{C}^r\) . In this number we give \(\mathcal{S}^r(\mathrm{X};\mathrm{E})\) the topology induced by the topology of compact \(\mathrm{C}^r\) -convergence on \(\mathcal{C}^r(\mathrm{X};\mathrm{E})\) , also called the topology of compact \(\mathrm{C}^r\) -convergence; it makes \(\mathcal{S}^r (\mathrm{X};\mathrm{E})\) into a complete separated locally convex topological vector space (cf.~Differentiable and Analytic Manifolds, Results, 15.3.1 and Spectral Theories, in preparation).

Now let H be a Lie group, \(m: H \times X \to X\) a law of left operation of class \(C^{r}\) ; put \(hx = m(h, x)\) for \(h \in H\) , \(x \in X\) . Let E be a vector H-bundle with base X, of class \(C^{r}\) (Chap. III, §1, no. 8, Def. 4). For \(s \in \mathcal{S}^{r}(X; E)\) and \(h \in H\) , denote by \(^{h}s\) the section \(x \mapsto h.s(h^{-1}x)\) of E; the map \((h, s) \mapsto ^{h}s\) is a law of operation of H on the space \(\mathcal{S}^{r}(X; E)\) .

Lemma 4. The law of operation \(\mathrm{H} \times \mathcal{S}^r (\mathrm{X};\mathrm{E}) \to \mathcal{S}^r (\mathrm{X};\mathrm{E})\) is continuous.

In view of the definition of the topology of \(\mathcal{S}^{r}(\mathrm{X};\mathrm{E})\) and General Topology, Chap. X, §3, no. 4, Th. 3, it suffices to prove that for any integers \(k \leq r\) , the map \(f : H \times X \times \mathcal{S}^{k}(\mathrm{X};\mathrm{E}) \to \mathrm{J}^{k}(\mathrm{X};\mathrm{E})\) such that \(f(h,x,s) = j_{x}^{k}(^{h}s)\) is continuous. For \(h \in H\) , denote by \(\tau_{h}\) (resp. \(\theta_{h}\) ) the automorphism \(x \mapsto hx\) of X (resp. of E). Define maps

\[
f _ {1}: \mathrm{H} \times \mathrm{X} \to \mathrm{J} ^ {k} (\mathrm{X}, \mathrm{X})
\]

\[
f _ {2}: \mathrm{H} \times \mathrm{E} \rightarrow \mathrm{J} ^ {k} (\mathrm{E}, \mathrm{E})
\]

\[
g: \mathrm{H} \times \mathrm{X} \times \mathcal {S} ^ {k} (\mathrm{X}; \mathrm{E}) \to \mathrm{J} ^ {k} (\mathrm{X}, \mathrm{E})
\]

by \(f_{1}(h,x) = j_{x}^{k}(\tau_{h}),f_{2}(h,v) = j_{v}^{k}(\theta_{h}),g(h,x,s) = j_{hx}^{k}(s)\) . We have

\[
f (h, x, s) = f _ {2} (h, s (h ^ {- 1} x)) \circ g (h ^ {- 1}, x, s) \circ f _ {1} (h ^ {- 1}, x),
\]

and consequently it suffices, by Differentiable and Analytic Manifolds, Results, 12.3.6, to prove that \(f_{1}, f_{2}\) and \(g\) are continuous.

Now \(g\) is the composite map

\[
\begin{array}{c c c} \mathrm{H} \times \mathrm{X} \times \mathcal {S} ^ {k} (\mathrm{X}; \mathrm{E}) & \xrightarrow {(m , \mathrm{Id})} & \mathrm{X} \times \mathcal {S} ^ {k} (\mathrm{X}; \mathrm{E}) \\ & \xrightarrow {(\mathrm{Id} , j ^ {k})} & \mathrm{X} \times \mathcal {C} (\mathrm{X}; \mathrm{J} ^ {k} (\mathrm{X}, \mathrm{E})) \quad \xrightarrow {\varepsilon} \quad \mathrm{J} ^ {k} (\mathrm{X}, \mathrm{E}) \end{array}
\]

with \(\varepsilon(x, u) = u(x)\) ; the map \(\varepsilon\) being continuous (General Topology, Chap. X, §3, no. 4, Cor. 1 of Th. 3), \(g\) is continuous.

Let \((h_{0}, x_{0}) \in \mathrm{H} \times \mathrm{X}\) ; we shall prove that \(f_{1}\) is continuous at \((h_{0}, x_{0})\) . There exist charts \((\mathrm{U}, \psi, \mathrm{F})\) and \((\mathrm{V}, \chi, \mathrm{F}')\) of \(\mathrm{X}\) and an open subset \(\Omega\) of \(\mathrm{H}\) such that \(x_{0} \in \mathrm{U}, h_{0} \in \Omega\) and \(m(\Omega \times \mathrm{U}) \subset \mathrm{V}\) . By using the expression for \(\mathrm{J}^{k}(\mathrm{X}, \mathrm{X})\) in these charts, we are reduced to proving, for \(1 \leq l \leq k\) , the continuity at \((h_{0}, x_{0})\) of the map \((h, x) \mapsto \Delta_{x}^{l}(\tau_{h})\) from \(\Omega \times \mathrm{U}\) to \(\mathrm{P}_{l}(\mathrm{F}; \mathrm{F}')\) , with \(\Delta_{x}^{l}(\tau_{h})(v) = \frac{1}{l!}\mathrm{D}^{l}\tau_{h}(x).v\) for \(v \in \mathrm{F}\) (Differentiable and Analytic Manifolds, Results, 12.2). But \(\mathrm{D}^{l}\tau_{h}(x)\) is simply the \(l\) th partial derivative of \(m(h, x)\) with respect to \(x\) , which is continuous by hypothesis; consequently, \(f_{1}\) is continuous. The proof that \(f_{2}\) is continuous is similar, hence the lemma.

PROPOSITION 5. Assume that the group H is compact and denote by dh the Haar measure on H of total mass 1. Let s be a section of E of class \(C^{r}\) .

Denote by \(s^{\sharp}\) the vector integral \(\int_{\mathrm{H}}{}^{h}s\,dh\) . Then \(s^{\sharp}\) is a section of \(\mathrm{E}\) of class \(\mathrm{C}^r\) , invariant under \(\mathrm{H}\) ; for \(x \in \mathrm{X}\) , we have \(s^{\sharp}(x) = \int_{\mathrm{H}}hs(h^{-1}x)\,dh \in \mathrm{E}_x\) . The endomorphism \(s \mapsto s^{\sharp}\) of \(\mathcal{S}^r(\mathrm{X};\mathrm{E})\) is a projection onto the subspace of H-invariant sections.

Consider the map \(h \mapsto {}^{h}s\) from H to \(\mathcal{S}^{r}(\mathrm{X};\mathrm{E})\) ; it is continuous by Lemma 4. Since the space \(\mathcal{S}^{r}(\mathrm{X};\mathrm{E})\) is separated and complete, the integral \(s^{\sharp} = \int_{H}{}^{h}s\,dh\) belongs to \(\mathcal{S}^{r}(\mathrm{X};\mathrm{E})\) (Integration, Chap. III, §3, no. 3, Cor. 2). The linear map \(s \mapsto s(x)\) from \(\mathcal{S}^{r}(\mathrm{X};\mathrm{E})\) to \(E_{x}\) being continuous, we have \(s^{\sharp}(x) = \int_{H}{}^{h}s(x)\,dh\) for all \(x \in X\) . It is clear that \(s^{\sharp}\) is invariant under H; if s is an H-invariant section, we have \(s^{\sharp} = s\) , hence the last assertion.

COROLLARY 1. Let \(\mathrm{F}\) be a Banach space, \(\rho : \mathrm{H} \to \mathbf{GL}(\mathrm{F})\) an analytic linear representation, \(f \in \mathcal{C}^r(\mathrm{X};\mathrm{F})\) . For \(x \in \mathrm{X}\) , put

\[
f ^ {\sharp} (x) = \int_ {\mathrm{H}} \rho (h). f (h ^ {- 1} x) d h.
\]

Then \(f^{\sharp}\) is a morphism of class \(\mathbf{C}^r\) from \(\mathbf{X}\) to \(\mathbf{F}\) , compatible with the operations of \(\mathbf{H}\) ; for \(x \in \mathbf{X}\) , we have (with \(\tau_h\) denoting the automorphism \(x \mapsto hx\) of \(\mathbf{X}\) )

\[
d _ {x} f ^ {\sharp} = \int_ {\mathrm{H}} (\rho (h) \circ d _ {h ^ {- 1} x} f \circ \mathrm{T} _ {x} (\tau_ {h ^ {- 1}})) d h \in \mathscr {L} (\mathrm{T} _ {x} (\mathrm{X}); \mathrm{F}).\tag{14}
\]

The first assertion follows from the proposition applied to the bundle \(X \times F\) , equipped with the law of operation \((h; (x, f)) \mapsto (hx, \rho(h).f)\) . The second follows from Integration, Chap. III, §3, no. 2, Prop. 2, by applying to the vector integral \(f^{\sharp}\) the homomorphism \(d_{x} : \mathcal{C}^{r}(\mathrm{X};\mathrm{F}) \to \mathcal{L}(\mathrm{T}_{x}(\mathrm{X});\mathrm{F})\) which is continuous by the definition of the topology of compact \(C^{r}\) -convergence.

COROLLARY 2. Let F be a Banach space, \(f \in \mathcal{C}^{r}(X;F)\) ; put

\[
f ^ {\sharp} (x) = \int_ {\mathrm{H}} f (h x) d h
\]

for \(x \in \mathrm{X}\) . The function \(f^{\sharp}\) is of class \(\mathbf{C}^r\) , and \(f^{\sharp}(hx) = f^{\sharp}(x)\) for \(x \in \mathrm{X}\) , \(h \in \mathrm{H}\) .

COROLLARY 3. Let F be a Banach space, p an integer \(\geq 0\) , \(^{k}\Omega^{p}(\mathrm{X};\mathrm{F})\) the space of differential forms of degree p on X, with values in F, and of class \(C^{k}\) ( \(2 \leq k + 1 \leq r\) ). For \(\omega \in {}^{k}\Omega^{p}(\mathrm{X};\mathrm{F})\) , put \(\omega^{\sharp} = \int_{H} \tau_{h}^{*}\omega dh\) . Then the map \(\omega \mapsto \omega^{\sharp}\) is a projection on \(^{k}\Omega^{p}(\mathrm{X};\mathrm{F})\) whose image is the subspace of H-invariant forms. We have \(d(\omega^{\sharp}) = (d\omega)^{\sharp}\) for all \(\omega \in {}^{k}\Omega^{p}(\mathrm{X};\mathrm{F})\) .

The first assertion follows from the proposition applied to the vector H-bundle \(\operatorname{Alt}^{p}(\mathrm{T}(\mathrm{X});\mathrm{F})\) (Chap. III, §1, no. 8, Examples). To prove the second assertion, it suffices, in view of Integration, Chap. III, §3, no. 2, Prop. 2, to prove that the map \(d:{}^{k}\varOmega^{p}(\mathrm{X};\mathrm{F})\to{}^{k-1}\varOmega^{p+1}(\mathrm{X};\mathrm{F})\) is continuous when the first (resp. second) space is given the topology of compact \(C^{k}\) -convergence (resp. \(C^{k-1}\) -convergence). But this follows immediately from the definition of these topologies by means of semi-norms (Spectral Theories, in preparation) and the fact that d is a differential operator of order \(\leq 1\) (Differentiable and Analytic Manifolds, Results, 14.4.2).

\subsection{INVARIANT DIFFERENTIAL FORMS}

Let X be a locally finite dimensional real manifold of class \(C^{\infty}\) , and let \((g, x) \mapsto gx\) be a law of left operation of class \(C^{\infty}\) of a connected compact Lie group G on X. For \(g \in G\) , denote by \(\tau_{g}\) the automorphism \(x \mapsto gx\) of X. Denote by \(\Omega(X)\) the algebra of real differential forms of class \(C^{\infty}\) on X (Differentiable and Analytic Manifolds, Results, 8.3.1).

For any element \(\xi\) of \(\mathfrak{g}\) , denote by \(D_{\xi}\) the corresponding vector field on X (Chap. III, §3, no. 5) and by \(\theta(\xi), i(\xi)\) the corresponding operators on \(\Omega(X)\) , so that we have the formulas (Differentiable and Analytic Manifolds, Results, 8.4.5 and 8.4.7)

\[
\theta (\xi) \omega = d (i (\xi) \omega) + i (\xi) d \omega\tag{15}
\]

\[
\frac {d}{d t} (\tau_ {\mathrm{exp} t \xi} ^ {*} \omega) = \tau_ {\mathrm{exp} t \xi} ^ {*} (\theta (\xi) \omega).\tag{16}
\]

A differential form \(\omega \in \Omega(\mathrm{X})\) is invariant if \(\tau_g^*\omega = \omega\) for all \(g \in \mathrm{G}\) ; by formula (16), this is equivalent to \(\theta(\xi)\omega = 0\) for all \(\xi \in \mathfrak{g}\) . Denote by \(\Omega(\mathrm{X})^{\mathrm{G}}\) the space of invariant differential forms on X; if \(\omega \in \Omega(\mathrm{X})^{\mathrm{G}}\) , we have \(d\omega \in \Omega(\mathrm{X})^{\mathrm{G}}\) , so \(\Omega(\mathrm{X})^{\mathrm{G}}\) is a subcomplex of the complex \((\Omega(\mathrm{X}), d)\) .

THEOREM 2. The canonical injection \(\iota : \Omega(\mathrm{X})^{\mathrm{G}} \to \Omega(\mathrm{X})\) is a homotopism of complexes (Algèbre, Chap. X, p.~33, déf. 5); the map \(\omega \mapsto \omega^{\sharp} = \int_{\mathrm{G}} \tau_g^* \omega dg\) is a homotopism, inverse to it up to homotopy. In particular, the map \(\mathrm{H}(\iota) : \mathrm{H}(\Omega(\mathrm{X})^{\mathrm{G}}) \to \mathrm{H}(\Omega(\mathrm{X}))\) is bijective.

By Cor. 3 of no. 4, the map \(\omega \mapsto \omega^{\sharp}\) is a morphism of complexes from \(\Omega(\mathrm{X})\) to \(\Omega(\mathrm{X})^{\mathrm{G}}\) that induces the identity on the subcomplex \(\Omega(\mathrm{X})^{\mathrm{G}}\) ; thus, to prove the theorem it suffices to construct a homomorphism \(s: \Omega(\mathrm{X}) \to \Omega(\mathrm{X})\) , graded of degree \(-1\) , such that

\[
\omega^ {\sharp} - \omega = (d \circ s + s \circ d) (\omega) \quad \text { for   all } \omega \in \Omega (\mathrm{X}).\tag{17}
\]

By Lemma 1 of Integration, Chapter IX, §2, no. 4 and Remark 1 of §2, no. 2, there exists a positive measure \(d\xi\) on \(\mathfrak{g}\) of compact support whose image under the exponential map is equal to \(dg\) . For \(\omega \in \Omega(X)\) , put

\[
s (\omega) = \int_ {0} ^ {1} \left\{\int_ {\mathfrak {g}} \tau_ {\exp t \xi} ^ {*} (i (\xi) \omega). d \xi \right\} d t;
\]

we have to show that formula (17) is satisfied. As in the proof of Cor. 1 (no. 4), we verify the formula

\[
d s (\omega) = \int_ {0} ^ {1} \left\{\int_ {\mathfrak {g}} \tau_ {\exp t \xi} ^ {*} d (i (\xi) \omega). d \xi \right\} d t.
\]

We now deduce from formulas (15) and (16) the equalities

\[
\begin{array}{r l} & {d s (\omega) + s (d \omega) = \int_ {0} ^ {1} \left\{\int_ {\mathfrak {g}} \tau_ {\exp t \xi} ^ {*} (d (i (\xi) \omega) + i (\xi) d \omega). d \xi \right\} d t} \\ & {\qquad = \int_ {0} ^ {1} \left\{\int_ {\mathfrak {g}} \tau_ {\exp t \xi} ^ {*} (\theta (\xi) \omega). d \xi \right\} d t} \\ & {\qquad = \int_ {\mathfrak {g}} \left\{\int_ {0} ^ {1} \frac {d}{d t} (\tau_ {\exp t \xi} ^ {*} \omega) d t \right\} d \xi} \\ & {\qquad = \int_ {\mathfrak {g}} (\tau_ {\exp \xi} ^ {*} \omega - \omega) d \xi} \\ & {\qquad = \omega^ {\sharp} - \omega ,} \end{array}
\]

hence Th. 2.

We apply the theorem in the case X = G, for the action of G by left translations. Recall (Chap. III, §3, no. 13, Prop. 50) that associating to a differential form on G its value at the identity element gives an isomorphism from \(\Omega(\mathrm{G})^{\mathrm{G}}\) to the graded algebra \(\operatorname{Alt}(\mathfrak{g})\) of alternating multilinear forms on g. Identify \(\Omega(\mathrm{G})^{\mathrm{G}}\) with \(\operatorname{Alt}(\mathfrak{g})\) by means of this isomorphism. The operator d is then given by the formula (Chap. III, §3, no. 14, Prop. 51)

\[
\begin{array}{l} d \omega (a _ {1}, \ldots , a _ {p + 1}) \\ = \sum_ {i <   j} (- 1) ^ {i + j} \omega ([ a _ {i}, a _ {j} ], a _ {1}, \ldots , a _ {i - 1}, a _ {i + 1}, \ldots , a _ {j - 1}, a _ {j + 1}, \ldots , a _ {p + 1}) \end{array}
\]

for \(\omega\) in \(\mathrm{Alt}^p (\mathfrak{g})\) and \(a_1,\ldots ,a_{p + 1}\) in \(\mathfrak{g}\) .

For \(\xi \in \mathfrak{g}\) , let \(\mathrm{L}_{\xi}\) be the corresponding left-invariant vector field (defined by means of the action of G on itself by right translations, cf.~Chap. III, §3, no. 6). The operators \(\theta (\mathrm{L}_{\xi}), i(\mathrm{L}_{\xi})\) commute with the action of G on \(\Omega (\mathrm{G})\) defined by left translation, and hence induce operators \(\theta (\xi), i(\xi)\) on \(\Omega (\mathrm{G})^{\mathrm{G}}\) ; with the preceding identifications, these are expressed by the formulas (Differentiable and Analytic Manifolds, Results, 8.3.2 and 8.4.2)

\[
(\theta (\xi) \omega) (a _ {1}, \dots , a _ {p}) = - \sum_ {i} \omega (a _ {1}, \dots , a _ {i - 1}, [ \xi , a _ {i} ], a _ {i + 1}, \dots , a _ {p})
\]

\[
(i (\xi) \omega) (a _ {1}, \ldots , a _ {p - 1}) = \omega (\xi , a _ {1}, \ldots , a _ {p - 1})
\]

for \(\omega\) in \(\operatorname{Alt}^p (\mathfrak{g})\) and \(a_1,\ldots ,a_p\) in \(\mathfrak{g}\) .

The subcomplex \({}^{\mathrm{G}}\varOmega(\mathrm{G})^{\mathrm{G}}\) of biinvariant forms (Chap. III, §3, no. 13) can be identified with the subcomplex \(\operatorname{Alt}(\mathfrak{g})^{\mathrm{G}}\) of alternating multilinear forms on g invariant under the adjoint representation (that is, such that \(\theta(\xi)\omega=0\) for all \(\xi\in\mathfrak{g}\) ). Thus, we have a commutative diagram of complexes

\[
\begin{array}{c c c c c} ^ {\mathrm{G}} \varOmega (\mathrm{G}) ^ {\mathrm{G}} & \longrightarrow & \varOmega (\mathrm{G}) ^ {\mathrm{G}} & \longrightarrow & \varOmega (\mathrm{G}) \\ \Big \downarrow & & \Big \downarrow & & \\ \operatorname{Alt} (\mathfrak {g}) ^ {\mathrm{G}} & \longrightarrow & \operatorname{Alt} (\mathfrak {g}) & & \end{array}\tag{18}
\]

where the horizontal arrows are the canonical injections, and the vertical arrows are the isomorphisms induced by the map \(\omega\mapsto\omega(e)\) .

COROLLARY 1. a) In the diagram (18), all the morphisms are homotopisms. b) Let \(\omega \in \operatorname{Alt}(\mathfrak{g})\) . Then \(\omega\) belongs to \(\operatorname{Alt}(\mathfrak{g})^{\mathrm{G}}\) if and only if \(d\omega = 0\) and \(d(i(\xi)\omega) = 0\) for all \(\xi \in g\) . The differential of the complex \(\operatorname{Alt}(\mathfrak{g})^{\mathrm{G}}\) is zero. c) The graded vector space \(\mathrm{H}(\Omega(\mathrm{G}))\) is isomorphic to \(\operatorname{Alt}(\mathfrak{g})^{\mathrm{G}}\) .

The theorem, applied to the action of G on G by left translations (resp. to the action \(((g,h);x)\mapsto gxh^{-1}\) of G × G on G) implies that the canonical injection \(\Omega(\mathrm{G})^{\mathrm{G}}\to\Omega(\mathrm{G})\) (resp. \({}^{G}\Omega(\mathrm{G})^{\mathrm{G}}\to\Omega(\mathrm{G})\) ) is a homotopism; in view of Algèbre, Chap. X, p.~34, Cor., assertion a) follows.

We prove \(b)\) . By Prop. 51 of Chap. III, §3, no. 14, we have \(d\alpha = -d\alpha\) , that is \(d\alpha = 0\) , for every differential form \(\alpha\) on \(G\) that is left and right invariant. Thus, if \(\omega \in \mathrm{Alt}(\mathfrak{g})^G\) , then \(d\omega = 0\) , and consequently \(d(i(\xi)\omega) = \theta(\xi)\omega - i(\xi)d\omega = 0\) . Conversely, if \(d\omega = 0\) and \(d(i(\xi)d\omega) = 0\) , then \(\theta(\xi)\omega = 0\) . Assertion \(c)\) follows from \(a)\) and \(b)\) .

Remark. Consider the subcomplexes \(\mathrm{Z}(\varOmega(\mathrm{G}))\) and \(\mathrm{B}(\varOmega(\mathrm{G}))\) of \(\varOmega(\mathrm{G})\) (Algèbre, Chap. X, p.~25). It follows from the formula giving the differential of the product of two forms (Differentiable and Analytic Manifolds, Results, 8.3.5) that \(\mathrm{Z}(\varOmega(\mathrm{G}))\) is a subalgebra of \(\varOmega(\mathrm{G})\) and that \(\mathrm{B}(\varOmega(\mathrm{G}))\) is an ideal of \(\mathrm{Z}(\varOmega(\mathrm{G}))\) ; consequently, the exterior product induces a graded algebra structure on \(\mathrm{H}(\varOmega(\mathrm{G}))\) . The preceding now gives an isomorphism of graded algebras \(\mathrm{H}(\varOmega(\mathrm{G})) \to \mathrm{Alt}(\mathfrak{g})^{\mathrm{G}}\) .

Let H be a closed subgroup of G; we apply Th. 2 to X = G/H. By Chap. III, §1, no. 8, Cor. 1 of Prop. 17, the G-invariant differential forms on G/H can be identified with the H-invariant elements of \(\operatorname{Alt}(\mathrm{T}_{e}(\mathrm{G}/\mathrm{H}))\) , that is with the elements of \(\operatorname{Alt}(\mathfrak{g})\) that are H-invariant and annihilated by the operators \(i(\xi)\) for all \(\xi \in \mathrm{L}(\mathrm{H})\) . Consequently:

COROLLARY 2. Let \(\mathrm{H}\) be a closed subgroup of \(\mathrm{G}\) .

\begin{enumerate}
\def\labelenumi{\alph{enumi})}
\item
  The canonical injection \(\Omega (\mathrm{G / H})^{\mathrm{G}}\to \Omega (\mathrm{G / H})\) is a homotopism.
\item
  The complex \(\Omega(\mathrm{G}/\mathrm{H})^{\mathrm{G}}\) can be identified with the subcomplex of \(\operatorname{Alt}(\mathfrak{g})\) consisting of the elements \(\omega\) of \(\operatorname{Alt}(\mathfrak{g})\) that are invariant under the adjoint representation of H and are such that \(i(\xi)\omega = 0\) for all \(\xi \in \operatorname{L}(\operatorname{H})\) . If, in addition, H is connected, this subcomplex consists of the \(\omega \in \operatorname{Alt}(\mathfrak{g})\) such that \(\theta(\xi)\omega = 0\) and \(i(\xi)\omega = 0\) for all \(\xi \in \operatorname{L}(\operatorname{H})\) .
\end{enumerate}

\section{IRREDUCIBLE REPRESENTATIONS OF CONNECTED COMPACT LIE GROUPS}

We retain the notations of §6. A representation of G is a continuous (hence analytic) homomorphism from G to a group \(\mathbf{GL}(\mathrm{V})\) , where V is a finite dimensional complex vector space. Every representation of G is semi-simple (§1, no. 1).

Choose a chamber C of \(\mathfrak{t}\) ( \(\S 5\) , no. 2), and put \(\Gamma(\mathrm{T})_{++} = \overline{\mathrm{C}} \cap \Gamma(\mathrm{T})\) .

\subsection{DOMINANT CHARACTERS}

Denote by \(X_{+}\) the set of elements \(\lambda\) of \(\mathrm{X(T)}\) such that \(\langle\lambda,x\rangle\geq0\) for all \(x\in\Gamma(\mathrm{T})_{++}\) , that is, such that the linear form \(\delta(\lambda):\mathfrak{t}_{\mathbf{C}}\to\mathbf{C}\) maps the chamber C of t to \(iR_{+}\) .

Give X(T) the ordered group structure for which the positive elements are those of \(X_{+}\) ; put \(R_{+} = R(G, T) \cap X_{+}\) and \(R_{-} = -R_{+}\) . The elements of \(R_{+}\) are called positive roots, those of \(R_{-}\) negative roots; every root is either positive or negative (Chap. VI, §1, no. 6, Th. 3). A positive root that is not the sum of two positive roots is said to be simple; every positive root is a sum of simple roots (loc. cit.); the simple roots form a basis of the subgroup of X(T) generated by the roots, a subgroup that can be identified with X(T/C(G)) (§4, no. 4); the reflections with respect to the simple roots generate the Weyl group \(W = W_{G}(T)\) (Chap. VI, §1, no. 5, Th. 2).

Lemma 1. Let \(\lambda\) be an element of \(\mathrm{X(T)}\) . The following conditions are equivalent:

\begin{enumerate}
\def\labelenumi{(\roman{enumi})}
\item
  \(\lambda - w(\lambda) \geq 0\) (resp. \(>0\) ) for all \(w \in \mathrm{W}\) such that \(w \neq 1\) ;
\item
  for all \(w \in \mathrm{W}\) such that \(w \neq 1\) , \(\lambda - w(\lambda)\) is a linear combination with positive coefficients (resp. positive coefficients not all zero) of simple roots;
\item
  \(\langle \lambda, K_{\alpha} \rangle \geq 0\) (resp. \(>0\) ) for every positive root \(\alpha\) ;
\item
  \(\langle\lambda,K_{\alpha}\rangle\geq0\) (resp. >0) for every simple root \(\alpha\) .
\end{enumerate}

The equivalence if (iii) and (iv) is immediate. Since the set of the \(K_{\alpha}\) can be identified with the inverse root system of \(\mathrm{R}(\mathrm{G},\mathrm{T})\) (§4, no. 5), the equivalence of (i) and (iii) follows from Chap. VI, §1, no. 6, Prop. 18 and Cor. The implication (ii) \(\Rightarrow\) (i) is trivial, and the opposite implication follows from loc. cit.

Denote by \(X_{++}\) the set of elements of \(\mathrm{X(T)}\) such that \(\langle\lambda,K_{\alpha}\rangle\geq0\) for every positive root \(\alpha\) . The elements of \(X_{++}\) are said to be dominant. They form a fundamental domain for the operation of W on \(\mathrm{X(T)}\) (Chap. VI, §1, no. 10). We have \(X_{++}\subset X_{+}\) .

If G is simply-connected, for each simple root \(\alpha\) there exists an element \(\varpi_{\alpha}\) of X(T) such that \(\langle\varpi_{\beta}, K_{\alpha}\rangle = \delta_{\alpha\beta}\) for every simple root \(\beta\) , that is, \(s_{\alpha}(\varpi_{\alpha}) = \varpi_{\alpha} - \alpha\) , \(s_{\beta}(\varpi_{\alpha}) = \varpi_{\alpha}\) for every simple root \(\beta \neq \alpha\) ; the \(\varpi_{\alpha}\) are called the fundamental dominant weights; they form a basis of the commutative group \(\mathrm{X(T)}\) and of the commutative monoid \(\mathrm{X}_{++}\) ; more precisely, every element \(\lambda\) of \(\mathrm{X(T)}\) can be written in the form \(\lambda = \sum_{\alpha} \langle \lambda, K_{\alpha} \rangle \varpi_{\alpha}\) .

Denote by \(\rho\) the element of \(\mathrm{X(T)}\otimes \mathbf{Q}\) such that

\[
2 \rho = \sum_ {\alpha \in \mathrm{R} _ {+}} \alpha .
\]

We have \(\langle\rho,K_{\alpha}\rangle=1\) for every simple root \(\alpha\) (Chap. VI, §1, no. 10, Prop. 29). If G is simply-connected, \(\rho\) is the sum of the fundamental dominant weights.

\subsection{HIGHEST WEIGHT OF AN IRREDUCIBLE REPRESENTATION}

Associate to any representation \(\tau : \mathrm{G} \to \mathrm{GL}(\mathrm{V})\) the homomorphism \(\mathrm{L}(\tau)_{(\mathbf{C})}\) from the \(\mathbf{C}\) -Lie algebra \(\mathfrak{g}_{\mathbf{C}}\) to \(\operatorname{End}(\mathrm{V})\) extending the linear representation \(\mathrm{L}(\tau)\) of \(\mathfrak{g}\) on the real vector space underlying V (Chap. III, §3, no. 11). By Prop. 7 of §4, no. 3, the map \(\delta\) from \(\mathrm{X}(\mathrm{T})\) to \(\operatorname{Hom}_{\mathbf{C}}(\mathfrak{t}_{\mathbf{C}}, \mathbf{C}) = \mathfrak{t}_{\mathbf{C}}^{*}\) induces a bijection from the set of weights of \(\tau\) relative to T to the set of weights of \(\mathrm{L}(\tau)_{(\mathbf{C})}\) relative to the Cartan subalgebra \(\mathfrak{t}_{\mathbf{C}}\) of \(\mathfrak{g}_{\mathbf{C}}\) .

Lemma 2. Let \(\varphi\) be a linear representation of the complex Lie algebra \(\mathfrak{g}_{\mathbf{C}}\) on a finite dimensional complex vector space V. There exists a representation \(\tau\) of G on V such that \(\mathrm{L}(\tau)_{(\mathbf{C})} = \varphi\) if and only if \(\varphi\) is semi-simple and the weights of \(\mathbf{t}_{\mathbf{C}}\) on V belong to \(\delta (\mathrm{X(T)})\) .

If there exists a representation \(\tau\) of G such that \(\mathrm{L}(\tau)_{(\mathbf{C})} = \varphi\) , then \(\varphi\) is semi-simple because G is connected and \(\tau\) is semi-simple (Chap. III, §6, no. 5, Cor. 2 of Prop. 13), and the weights of \(\mathfrak{t}_{\mathbf{C}}\) on V belong to the image of \(\delta\) . Thus, the condition is necessary; we shall prove that it is sufficient. If \(\varphi\) is semi-simple, V is the direct sum of the \(V_{\mu}(t_{\mathbf{C}})\) , where \(\mu\) belongs to the set of weights of \(\mathfrak{t}_{\mathbf{C}}\) on V (Chap. VII, §2, no. 4, Cor. 3 of Th. 2); if all the weights belong to the image of \(\delta\) , there exists a representation \(\tau_{\mathrm{T}}\) of T on V such that \(\mathrm{L}(\tau_{\mathrm{T}})_{(\mathbf{C})} = \varphi |t_{\mathbf{C}}\) : indeed, it suffices to put \(\tau_{\mathrm{T}}(t)v = t^{\lambda}v\) for \(t \in \mathrm{T}\) and \(v \in V_{\delta(\lambda)}(t_{\mathbf{C}})\) . The lemma now follows from Prop. 8 of §2, no. 6.

THEOREM 1. a) Let \(\tau : \mathbf{G} \to \mathbf{GL}(\mathbf{V})\) be an irreducible representation of \(\mathbf{G}\) . Then the set of weights of \(\tau\) (relative to \(\mathbf{T}\) ) has a largest element \(\lambda\) , which is dominant, and the space \(\mathrm{V}_{\lambda}(\mathrm{T})\) is of dimension 1.

\begin{enumerate}
\def\labelenumi{\alph{enumi})}
\setcounter{enumi}{1}
\item
  Two irreducible representations of G are equivalent if and only if their highest weights are equal.
\item
  For every dominant element \(\lambda\) of X(T), there exists an irreducible representation of G of highest weight \(\lambda\) .
\end{enumerate}

By Lemma 2, the equivalence classes of irreducible representations of G correspond bijectively to the classes of finite dimensional irreducible representations of g whose weights belong to \(\delta(\mathrm{X}(\mathrm{T}))\) .

Denote by \(\mathcal{G}\mathfrak{g}_{\mathbf{C}}\) the centre and by \(\mathcal{D}\mathfrak{g}_{\mathbf{C}}\) the derived Lie algebra of \(\mathfrak{g}_{\mathbf{C}}\) , so that \(\mathfrak{g}_{\mathbf{C}} = \mathcal{G}\mathfrak{g}_{\mathbf{C}} \oplus \mathcal{D}\mathfrak{g}_{\mathbf{C}}\) . For every bilinear form \(\mu\) on \(\mathfrak{t}_{\mathbf{C}} \cap \mathcal{D}\mathfrak{g}_{\mathbf{C}}\) , denote by \(\mathrm{E}(\mu)\) the simple \(\mathcal{D}\mathfrak{g}_{\mathbf{C}}\) -module introduced in Chap. VIII, §6, no. 3; for every linear form \(\nu\) on \(\mathcal{G}\mathfrak{g}_{\mathbf{C}}\) , denote by \(\mathbf{C}(\nu)\) the 1-dimensional \(\mathcal{G}\mathfrak{g}_{\mathbf{C}}\) -module over \(\mathbf{C}\) associated to it. Then the \(\mathfrak{g}_{\mathbf{C}}\) -modules \(\mathbf{C}(\nu) \otimes \mathrm{E}(\mu)\) are simple, and by Chap. VIII, §7, no. 2, Cor. 2 of Th. 1 and Algebra, Chap. VIII, §11, no. 1, Th. 1, every finite dimensional simple \(\mathfrak{g}_{\mathbf{C}}\) -module is isomorphic to one of the modules \(\mathbf{C}(\nu) \otimes \mathrm{E}(\mu)\) ; moreover (loc. cit.) \(\mathbf{C}(\nu) \otimes \mathrm{E}(\mu)\) is finite dimensional if and only if \(\mu(H_{\alpha})\) is a positive integer for every simple root \(\alpha\) . If we denote by \(\nu + \mu\) the linear form on \(\mathfrak{t}_{\mathbf{C}}\) that induces \(\nu\) on \(\mathcal{G}\mathfrak{g}_{\mathbf{C}}\) and \(\mu\) on \(\mathfrak{t}_{\mathbf{C}} \cap \mathcal{D}\mathfrak{g}_{\mathbf{C}}\) , then \((\nu + \mu)(H_{\alpha}) = \mu(H_{\alpha})\) ; moreover, the weights of \(\mathbf{C}(\nu) \otimes \mathrm{E}(\mu)\) are the \(\nu + \lambda\) , where \(\lambda\) is any weight of \(\mathrm{E}(\mu)\) , and are thus of the form \(\nu + \mu - \theta\) , with \(\theta \in \delta(X_{+})\) (Chap. VIII, §6, no. 2, Lemma 2).

We conclude that the g-module \(\mathbf{C}(\nu)\otimes\mathrm{E}(\mu)\) is finite dimensional if and only if \((\nu+\mu)(H_{\alpha})\) is a positive integer for every simple root \(\alpha\) , and that its weights belong to \(\delta(\mathrm{X}(\mathrm{T}))\) if and only if \(\nu+\mu\) belongs to \(\delta(\mathrm{X}(\mathrm{T}))\) . The conjunction of these two conditions means that \(\nu+\mu\) belongs to \(\delta(\mathrm{X}_{++})\) ; in that case, \(\nu+\mu\) is the highest weight of \(\mathbf{C}(\nu)\otimes\mathrm{E}(\mu)\) . Thus, we have constructed for every dominant weight \(\lambda\) of X(T) an irreducible representation of G of highest weight \(\lambda\) , and have thus obtained, up to equivalence, all the irreducible representations of G. Since the vectors of weight \(\nu+\mu\) in \(\mathbf{C}(\nu)\otimes\mathrm{E}(\mu)\) form a subspace of dimension 1, the proof is complete.

COROLLARY. The group G has a (finite dimensional) faithful linear representation.

Observe first that every element of X(T) is equal to the difference of two dominant weights: more precisely, let \(\varpi\) be an element of \(X_{++}\) such that \(\langle\varpi,K_{\alpha}\rangle>0\) for every simple root \(\alpha\) ; for all \(\lambda\in X(T)\) there exists a positive integer n such that \(\langle\lambda+n\varpi,K_{\alpha}\rangle\geq0\) for every simple root \(\alpha\) , that is (no. 1, Lemma 1) \(\lambda+n\varpi\in X_{++}\) .

Consequently, there exists a finite family \((\lambda_{i})_{i\in I}\) of elements of \(X_{++}\) generating the Z-module \(\mathrm{X(T)}\) . For \(i\in I\) , let \(\tau_{i}\) be an irreducible representation of G of highest weight \(\lambda_{i}\) (Th. 1); let the representation \(\tau\) be the direct sum of the \(\tau_{i}\) . By construction the set \(\mathrm{P}(\tau,\mathrm{T})\) of weights of \(\tau\) (relative to T) generates the Z-module \(\mathrm{X(T)}\) . It now follows from Prop. 6 of §4, no. 3 that the homomorphism \(\tau\) is injective, hence the corollary.

Remarks. 1) Let \(\mathfrak{n}_{+}\) be the subalgebra of \(\mathfrak{g}_{\mathbf{C}}\) that is the sum of the \(\mathfrak{g}^{\alpha}\) for \(\alpha > 0\) . Let \(\tau: \mathrm{G} \to \mathrm{GL}(\mathrm{V})\) be an irreducible representation, \(\lambda \in \mathrm{X}_{++}\) its highest weight and \(\tau': \mathfrak{g}_{\mathbf{C}} \to \mathfrak{gl}(\mathrm{V})\) the representation induced by \(\tau\) . Then \(\mathrm{V}_{\lambda}(\mathrm{T})\) is the subspace consisting of the vectors \(v\) in \(\mathrm{V}\) such that \(\tau'(x)v = 0\) for all \(x \in \mathfrak{n}_{+}\) .

Indeed, this follows from the corresponding statement for the \(\mathfrak{g}_{\mathbf{C}}\) -modules \(\mathbf{C}(\nu) \otimes \mathrm{E}(\mu)\) (Chap. VIII, §6, no. 2, Prop. 3).

\begin{enumerate}
\def\labelenumi{\arabic{enumi})}
\setcounter{enumi}{1}
\item
  Let \(\Theta(\mathrm{G})\) be the algebra of continuous representative functions on \(\mathrm{G}\) with values in \(\mathbf{C}\) (Algebra, Chap. VIII). Let \(\mathrm{G}\) operate on \(\Theta(\mathrm{G})\) by left and right translations. For each \(\lambda \in \mathrm{X}_{++}\) , let \((\mathrm{V}_{\lambda}, \tau_{\lambda})\) be an irreducible representation of \(\mathrm{G}\) of highest weight \(\lambda\) (Th. 1), and \((\mathrm{V}_{\lambda}^{*}, \check{\tau}_{\lambda})\) the contragredient representation (Chap. III, §3, no. 11); by Spectral Theories, the representation of \(\mathrm{G} \times \mathrm{G}\) on \(\Theta(\mathrm{G})\) is isomorphic to the direct sum of the representations \((\mathrm{V}_{\lambda} \otimes \mathrm{V}_{\lambda}^{*}, \tau_{\lambda} \otimes \check{\tau}_{\lambda})\) for all \(\lambda\) in \(\mathrm{X}_{++}\) . Remark 1 now implies the following statement: Let \(\lambda \in \mathrm{X}_{++}\) , and let \(\mathrm{E}_{\lambda}\) be the vector subspace of \(\Theta(\mathrm{G})\) consisting of the continuous representative functions \(f\) on \(\mathrm{G}\) such that \(f(gt) = \lambda(t)^{-1} f(g)\) for all \(g \in \mathrm{G}\) and all \(t \in \mathrm{T}\) , and such that \(f * x = 0\) for all \(x \in \mathfrak{n}_{-} = \bigoplus_{\alpha < 0} \mathfrak{g}^{\alpha}\) . Then \(\mathrm{E}_{\lambda}\) is stable under left translations, and the representation of \(\mathrm{G}\) on \(\mathrm{E}_{\lambda}\) by left translations is irreducible of highest weight \(\lambda\) .
\item
  Let \(\tau : \mathrm{G} \to \mathbf{GL}(\mathrm{V})\) be an irreducible representation. There exists an element \(\nu\) of \(\mathrm{X}(\mathrm{C}(\mathrm{G}))\) such that \(\tau(s)v = \nu(s)v\) for all \(s \in \mathrm{C}(\mathrm{G}), v \in \mathrm{V}\) : indeed, \(\tau(\mathrm{C}(\mathrm{G}))\) is contained in the commutant of \(\tau(\mathrm{G})\) , which is equal to \(\mathbf{C}^*.1_{\mathrm{V}}\) (Algebra, Chap. VIII, §3, no. 2, Th. 1). For every weight \(\lambda\) of \(\tau\) , the restriction of \(\lambda\) to \(\mathrm{C}(\mathrm{G})\) is equal to \(\nu\) .
\item
  The definitions and statements of Chap. VIII, §7, nos. 2 to 5 can be generalized without difficulty to the present situation; we leave the details to the reader.
\end{enumerate}

PROPOSITION 1. Let \(\tau : G \to \mathbf{GL}(V)\) be an irreducible representation of \(G\) of highest weight \(\lambda \in X_{++}\) . Let \(m\) be the integer \(\sum_{\alpha \in R_+} \langle \lambda, K_\alpha \rangle\) , and let \(w_0\) be the element of the Weyl group such that \(w_0(R_+) = R_-\) (Chap. VI, § 6, Cor. 3 of Prop. 17). There are three possible cases:

\begin{enumerate}
\def\labelenumi{\alph{enumi})}
\item
  \(w_{0}(\lambda) = -\lambda\) and m is even. Then there exists a G-invariant separating symmetric bilinear form on V; the representation \(\tau\) is of real type (Appendix II).
\item
  \(w_{0}(\lambda) \neq -\lambda\) . Every G-invariant bilinear form on V is zero; the representation \(\tau\) is of complex type (loc. cit.).
\item
  \(w_{0}(\lambda) = -\lambda\) and m is odd. Then there exists a G-invariant separating alternating bilinear form on V; the representation \(\tau\) is of quaternionic type (loc. cit.).
\end{enumerate}

If the restriction of \(\tau\) to \(\mathrm{C}(\mathrm{G})_0\) is not trivial, we are in case \(b\) ).

A bilinear form B on V is invariant under G if and only if it is invariant under \(g_{C}\) (Chap. III, §6, no. 5, Cor. 3). Thus, if G is semi-simple, the proposition follows from Chap. VIII, §7, no. 5, Prop. 12 and Prop. 3 of Appendix II.

In the general case, put \(\mathrm{C}(\mathrm{G})_{0} = \mathrm{S}\) , and identify \(\mathrm{X}(\mathrm{T}/\mathrm{S})\) with a subgroup of \(\mathrm{X}(\mathrm{T})\) (stable under W). If \(\tau(\mathrm{S}) = \{1_{\mathrm{V}}\}\) , \(\tau\) induces by passage to the quotient a representation \(\tau' : G/S \to GL(V)\) of highest weight \(\lambda\) ; in this case the proposition follows from the preceding, applied to G/S.

Assume that \(\tau (\mathrm{S})\neq \{1_{\mathrm{V}}\}\) . There exists a non-zero element \(\nu\) of \(\mathrm{X(S)}\) such that \(\tau (s) = \nu (s)_{\mathrm{V}}\) for all \(s\in \mathrm{S}\) (Remark 3). Then \(\nu\) is the image of \(\lambda\) under the restriction homomorphism \(\mathrm{X}(\mathrm{T})\to\mathrm{X}(\mathrm{S})\) ; since W operates trivially on \(\mathrm{X}(\mathrm{S})\) , the equality \(w_{0}(\lambda)=-\lambda\) implies that \(\nu=-\nu\) , which is impossible: hence \(w_{0}(\lambda)\neq-\lambda\) . On the other hand, if B is a G-invariant bilinear form on V, we have, for all x,y in V and s in S,

\[
\mathrm{B} (\nu (s) x, \nu (s) y) = \mathrm{B} (x, y) = \nu (s) ^ {2} \mathrm{B} (x, y)
\]

which implies that B = 0, hence the proposition.

Let \(\mathfrak{S}_{\mathbf{R}}(\mathrm{G})\) be the set of classes of irreducible continuous representations of G on finite dimensional real vector spaces. Prop. 1 and the results of Appendix II give a bijection \(\Phi:X_{++}/\Sigma\to\mathfrak{S}_{\mathbf{R}}(\mathrm{G})\) , where \(\Sigma\) denotes the subgroup \(\{1,-w_{0}\}\) of \(\operatorname{Aut}(\mathrm{X}(\mathrm{T}))\) . More precisely, let \(\lambda\in X_{++}\) , and let \(E_{\lambda}\) be a representation of G of highest weight \(\lambda\) ; then

\[
\begin{array}{c} \Phi (\{\lambda , - w _ {0} (\lambda) \}) = \mathrm{E} _ {\lambda [ \mathbf {R} ]} \text {if} \lambda \neq - w _ {0} (\lambda) \text {or if} \sum_ {\alpha \in \mathrm{R} _ {+}} \langle \lambda , K _ {\alpha} \rangle \notin 2 \mathbf {Z} \\ \Phi (\{\lambda , - w _ {0} (\lambda) \}) = \mathrm{E} _ {\lambda} ^ {\prime} \text {if} \lambda = - w _ {0} (\lambda) \text {and} \sum_ {\alpha \in \mathrm{R} _ {+}} \langle \lambda , K _ {\alpha} \rangle \in 2 \mathbf {Z} \end{array}
\]

where \(E_{\lambda}^{\prime}\) is an R-structure on \(E_{\lambda}\) invariant under G.

\subsection{THE RING R(G)}

Let \(\mathrm{R}(\mathrm{G})\) be the ring of representations (continuous, on finite dimensional complex vector spaces) of \(\mathrm{G}\) (Algebra, Chap. VIII, §10, no. 6). If \(\tau\) is a representation of \(\mathrm{G}\) , denote by \([\tau]\) its class in \(\mathrm{R}(\mathrm{G})\) ; if \(\tau\) and \(\tau'\) are two representations of \(\mathrm{G}\) , we have, by definition,

\[
\begin{array}{r} [ \tau ] + [ \tau^ {\prime} ] = [ \tau \oplus \tau^ {\prime} ] \\ \relax [ \tau ] [ \tau^ {\prime} ] = [ \tau \otimes \tau^ {\prime} ]. \end{array}
\]

Since every representation of G is semi-simple, the Z-module \(\mathrm{R}(\mathrm{G})\) is free and admits as a basis the set of classes of irreducible representations of G, a set that can be identified with \(X_{++}\) by Th. 1. The map \(\tau \mapsto \mathrm{L}(\tau)_{(\mathbf{C})}\) induces a homomorphism l from the ring \(\mathrm{R}(\mathrm{G})\) to the ring \(\mathcal{R}(\mathfrak{g}_{\mathbf{C}})\) of representations of \(g_{C}\) (Chap. VIII, §7, no. 6).

Let \(\tau : \mathbf{G} \to \mathbf{GL}(\mathbf{V})\) be a representation of \(\mathbf{G}\) ; we consider the gradation \((\mathrm{V}_{\lambda}(\mathrm{T}))_{\lambda \in \mathrm{X}(\mathrm{T})}\) of the \(\mathbf{C}\) -vector space \(\mathbf{V}\) . Denote by \(\mathrm{Ch}(\mathbf{V})\) , or by \(\mathrm{Ch}(\tau)\) , the character of the graded vector space \(\mathbf{V}\) (Chap. VIII, §7, no. 7); if \((e^{\lambda})_{\lambda \in \mathrm{X}(\mathrm{T})}\) denotes the canonical basis of the algebra \(\mathbf{Z}[\mathrm{X}(\mathrm{T})] = \mathbf{Z}^{(\mathrm{X}(\mathrm{T}))}\) , we have, by definition,

\[
\operatorname{Ch} (\tau) = \sum_ {\lambda \in \mathrm{X} (\mathrm{T})} \left(\dim \mathrm{V} _ {\lambda} (\mathrm{T})\right) e ^ {\lambda}.
\]

We define in the same way (loc. cit.) a homomorphism of rings, also denoted by Ch, from R(G) to Z[X(T)]. If G is semi-simple, we have a commutative diagram

\[
\begin{array}{c c c} \mathrm{R(G)} & \stackrel {{\mathrm{Ch}}} {{\longrightarrow}} & \mathbf {Z} [ \mathrm{X(T)} ] \\ \Big \downarrow_ {l} & & \Big \downarrow_ {\tilde {\delta}} \\ \mathcal {R} (\mathfrak {g} _ {\mathbf {C}}) & \stackrel {{\mathrm{ch}}} {{\longrightarrow}} & \mathbf {Z} [ \mathrm{P} ] \end{array}\tag{1}
\]

where P denotes the group of weights of \(\mathrm{R}(\mathfrak{g}_{\mathbf{C}}, \mathfrak{t}_{\mathbf{C}})\) and \(\tilde{\delta}\) the homomorphism induced by \(\delta\) .

The Weyl group W operates by automorphisms on the group X(T), and hence on the ring Z[X(T)]. By Prop. 5 of §4, no. 3, the image of Ch is contained in the subring Z[X(T)] \(^{W}\) consisting of the elements invariant under W.

PROPOSITION 2. The homomorphism \(\mathrm{Ch}\) induces an isomorphism from \(\mathrm{R}(\mathrm{G})\) to \(\mathrm{Z}[\mathrm{X}(\mathrm{T})]^{\mathrm{W}}\) .

For \(\lambda \in \mathrm{X}_{++}\) , denote by \([\lambda]\) the class in \(\mathrm{R}(\mathrm{G})\) of the irreducible representation of highest weight \(\lambda\) . Since the family \(([\lambda])_{\lambda \in \mathrm{X}_{++}}\) is a basis of the \(\mathbf{Z}\) -module \(\mathrm{R}(\mathrm{G})\) , it suffices to prove the following assertion:

The family \((\mathrm{Ch}[\lambda])_{\lambda\in\mathrm{X}_{++}}\) is a basis of the Z-module \(\mathbf{Z}[\mathrm{X}(\mathrm{T})]^{\mathrm{W}}\) .

For every element \(u = \sum_{\lambda} a_{\lambda} e^{\lambda}\) of \(\mathbf{Z}[X(T)]\) , a term \(a_{\lambda} e^{\lambda}\) in u is called maximal if \(\lambda\) is a maximal element of the set of \(\mu \in X(T)\) such that \(a_{\mu} \neq 0\) . Th. 1 implies that \(\operatorname{Ch}[\lambda]\) has a unique maximal term, namely \(e^{\lambda}\) . Thus, the proposition follows from the following lemma.

Lemma 3. For each \(\lambda \in \mathrm{X}_{++}\) , let \(\mathrm{C}_{\lambda}\) be an element of \(\mathbf{Z}[\mathrm{X(T)}]^{\mathrm{W}}\) having unique maximal term \(e^{\lambda}\) . Then the family \((\mathrm{C}_{\lambda})_{\lambda \in \mathrm{X}_{++}}\) is a basis of \(\mathbf{Z}[\mathrm{X(T)}]^{\mathrm{W}}\) .

The proof is identical to that of Prop. 3 of Chap. VI, §3, no. 4, replacing A by Z, P by X(T) and P ∩ \(\overline{C}\) by \(X_{++}\) .

Let \(\Theta (\mathrm{G})\) (resp. \(\Theta (\mathrm{T})\) ) be the \(\mathbf{C}\) -algebra of continuous representative functions on \(\mathrm{G}\) (resp. T), and let \(\mathrm{Z}\Theta (\mathrm{G})\) (resp. \(\Theta (\mathrm{T})^{\mathrm{W}}\) ) be the subalgebra consisting of the central (resp. W-invariant) functions. The restriction map \(\Theta (\mathrm{G})\to \Theta (\mathrm{T})\) induces a homomorphism of rings \(r:\mathrm{Z}\Theta (\mathrm{G})\rightarrow \Theta (\mathrm{T})^{\mathrm{W}}\) . On the other hand, the map that associates to a representation \(\tau\) its character (that is, the function \(g\mapsto \mathrm{Tr}\tau (g)\) ) extends to a homomorphism of \(\mathbf{C}\) -algebras \(\mathrm{Tr}:\mathbf{C}\otimes_{\mathbf{Z}}\mathrm{R}(\mathrm{G})\to \mathrm{Z}\Theta (\mathrm{G})\) which, by Spectral Theories, is an isomorphism. Similarly, the canonical injection \(\mathrm{X(T)}\to \Theta (\mathrm{T})\) induces an isomorphism of \(\mathbf{C}\) -algebras \(\iota :\mathbf{C}[\mathrm{X(T)}]\to \Theta (\mathrm{T})\) , which induces an isomorphism \(\iota :\mathbf{C}[\mathrm{X(T)}]^{\mathrm{W}}\to \Theta (\mathrm{T})^{\mathrm{W}}\) . The diagram

\[
\begin{array}{c c c} \mathbf {C} \otimes_ {\mathbf {Z}} \mathrm{R(G)} & \stackrel {{1 \otimes \mathrm{Ch}}} {{\longrightarrow}} & \mathbf {C} [ \mathrm{X(T)} ] ^ {\mathrm{W}} \\ \Big \downarrow_ {\mathrm{Tr}} & & \Big \downarrow_ {\iota} \\ \mathrm{Z} \Theta (\mathrm{G}) & \stackrel {{r}} {{\longrightarrow}} & \Theta (\mathrm{T}) ^ {\mathrm{W}} \end{array}\tag{2}
\]

is commutative: indeed, for every representation \(\tau : \mathrm{G} \to \mathbf{GL}(\mathrm{V})\) and all \(t \in \mathrm{T}\) ,

\[
\operatorname{Tr} \tau (t) = \sum_ {\lambda \in \mathrm{X} (\mathrm{T})} \left(\dim \mathrm{V} _ {\lambda} (\mathrm{T})\right) \lambda (t) = \iota (\operatorname{Ch} \tau) (t),
\]

that is, \((r\circ \mathrm{Tr})[\tau ] = (\iota \circ \mathrm{Ch})[\tau ]\)

We can now deduce from the proposition the following result.

COROLLARY. The restriction map \(r: \mathrm{Z}\Theta(\mathrm{G}) \to \Theta(\mathrm{T})^{\mathrm{W}}\) is bijective.

\subsection{CHARACTER FORMULA}

In this number, we write the group \(X(T)\) multiplicatively, and regard its elements as complex-valued functions on T. We assume that the element \(\rho\) of \(X(T) \otimes Q\) belongs to \(X(T)\) .

Denote by \(L^{2}(T)\) the Hilbert space of classes of square-integrable complex functions on T, and by \(\Theta(T)\) the subspace consisting of the continuous representative functions. By Spectral Theories, \(X(T)\) is an orthonormal basis of \(L^{2}(T)\) and an algebraic basis of \(\Theta(T)\) .

For \(f \in \mathrm{L}^{2}(\mathrm{T})\) and \(w \in W\) , denote by wf the element of \(\mathrm{L}^{2}(\mathrm{T})\) defined by \(wf(t) = f(w^{-1}(t))\) ; thus, for \(\lambda \in \mathrm{X}(\mathrm{T})\) we have \(w\lambda = w(\lambda)\) . Denote by \(\varepsilon : W \to \{1, -1\}\) the signature (the unique homomorphism such that \(\varepsilon(s) = -1\) for every reflection s); for \(f \in \mathrm{L}^{2}(\mathrm{T})\) , put

\[
\mathrm{J} (f) = \sum_ {w \in \mathrm{W}} \varepsilon (w) ^ {w} f.
\]

If \(\lambda \in X_{++}\) , the characters \(^w (\lambda \rho)\) are distinct; indeed, it suffices to prove that \(^w (\lambda \rho)\neq \lambda \rho\) for all \(w\neq 1\) ; but this follows from Lemma 1 (no. 1) and the fact that \(\langle \lambda \rho ,K_{\alpha}\rangle = \langle \lambda ,K_{\alpha}\rangle +1 > 0\) for every positive root \(\alpha\) . Consequently,

\[
\| \mathrm{J} (\lambda \rho) \| ^ {2} = \operatorname{Card} (\mathrm{W}) = w (\mathrm{G}).
\]

An element \(f \in \mathrm{L}^{2}(\mathrm{T})\) is said to be anti-invariant if \(^{w}f = \varepsilon(w)f\) for all \(w \in W\) (that is, if \(^{s}f = -f\) for every reflection s). We shall show that \(\frac{1}{w(\mathrm{G})}\mathrm{J}\) is the orthogonal projection of \(\mathrm{L}^{2}(\mathrm{T})\) onto the subspace of anti-invariant elements. Indeed, let \(f, f'\) be in \(\mathrm{L}^{2}(\mathrm{T})\) , with \(f'\) anti-invariant; then \(\mathrm{J}(f)\) is anti-invariant and

\[
\begin{array}{c} \langle f ^ {\prime}, \mathrm{J} (f) \rangle = \sum_ {w \in \mathrm{W}} \varepsilon (w) \langle f ^ {\prime},   ^ {w} f \rangle = \sum_ {w \in \mathrm{W}} \langle^ {w} f ^ {\prime},   ^ {w} f \rangle \\ = \sum_ {w \in \mathrm{W}} \langle f ^ {\prime},   f \rangle = w (\mathrm{G}) \langle f ^ {\prime},   f \rangle . \end{array}
\]

PROPOSITION 3. The elements \(\mathrm{J}(\lambda\rho)/\sqrt{w(\mathrm{G})}\) , for \(\lambda\in X_{++}\) , form an orthonormal basis of the subspace of anti-invariant elements of \(\mathrm{L}^{2}(\mathrm{T})\) , and an algebraic basis of the subspace of anti-invariant elements of \(\Theta(\mathrm{T})\) .

The proof is identical to that of Chap. VI, §3, no. 3, Prop. 1.

By Chap. VI, loc. cit., Prop. 2,

\[
\mathrm{J} (\rho) = \rho \prod_ {\alpha > 0} (1 - \alpha^ {- 1}) = \rho^ {- 1} \prod_ {\alpha > 0} (\alpha - 1),\tag{3}
\]

so

\[
\mathrm{J} (\rho) \overline {{\mathrm{J} (\rho)}} = \prod_ {\alpha} (\alpha - 1).\tag{4}
\]

By Cor. 2 of Th. 1 (§6, no. 2), we deduce:

Lemma 4. If \(\varphi\) and \(\psi\) are two continuous central functions on \(\mathbf{G}\) ,

\[
\int_ {\mathrm{G}} \overline {{\varphi (g)}} \psi (g) d g = \frac {1}{w (\mathrm{G})} \int_ {\mathrm{T}} \overline {{(\varphi (t) \mathrm{J} (\rho) (t))}}. (\psi (t) \mathrm{J} (\rho) (t)) d t.
\]

For all \(\lambda\in X_{++}\) , denote by \(\chi_{\lambda}\) the character of an irreducible representation of G of highest weight \(\lambda\) .

THEOREM 2 (H. Weyl). For all \(\lambda \in \mathrm{X}_{++}\) , we have \(\mathrm{J}(\rho)\chi_{\lambda}|\mathrm{T} = \mathrm{J}(\lambda \rho)\) .

The function \(\mathrm{J}(\rho)\cdot\chi_{\lambda}|T\) is anti-invariant under W, and is a linear combination with integer coefficients of elements of \(\mathrm{X(T)}\) . Thus, by Chap. VI, §3, no. 3, Prop. 1, it can be written as \(\sum_{\mu}a_{\mu}\mathrm{J}(\mu\rho)\) , where \(\mu\) belongs to \(X_{++}\) , and the \(a_{\mu}\) are integers all but finitely-many of which are zero; since \(\int_{\mathrm{G}}|\chi_{\lambda}(g)|^{2}dg=1\)

(Spectral Theories), it follows from Prop. 3 and Lemma 4 that \(\sum_{\mu}(a_{\mu})^{2} = 1\) ; thus, the \(a_{\mu}\) are all zero, except one which must be equal to 1 or \(-1\) . But the coefficient of \(\lambda\) in \(\chi_{\alpha}|T\) is equal to 1 (Th. 1), hence the coefficient of \(\lambda \rho\) in \(J(\rho)\cdot \chi_{\lambda}|T\) is equal to 1 (Chap. VI, §3, no. 3, Remark 2), which implies that \(a_{\lambda \rho} = 1\) , hence the theorem.

COROLLARY 1. With the notations of no. 3, we have in \(\mathbf{Z}[\mathrm{X(T)}]\) ,

\[
\left(\sum_ {w \in \mathrm{W}} \varepsilon (w) e ^ {w \rho}\right) \operatorname{Ch} [ \lambda ] = \sum_ {w \in \mathrm{W}} \varepsilon (w) e ^ {w \lambda} e ^ {w \rho} \text {for all} \lambda \in \mathrm{X} _ {+ +}.
\]

This follows from the theorem and commutative diagram (2) (no. 3).

COROLLARY 2. For all \(\lambda \in X_{++}\) and every regular element \(t\) of T,

\[
\chi_ {\lambda} (t) = \frac {\sum_ {w} \varepsilon (w) \lambda (w t) \rho (w t)}{\sum_ {w} \varepsilon (w) \rho (w t)}\tag{5}
\]

where the two sums are over the elements w of W.

Indeed, \(\mathrm{J}(\rho)(t)\) is non-zero, since \(t\) is regular (formula (4)).

If \(\varphi\) is a central function on \(\mathbf{G}\) , the restriction of \(\varphi\) to \(\mathrm{T}\) is invariant under \(\mathrm{W}\) , so \(\mathrm{J}(\rho). \varphi|\mathrm{T}\) is anti-invariant under \(\mathrm{W}\) . Further, by Spectral Theories and Th. 1, the family \((\chi_{\lambda})_{\lambda \in \mathrm{X}_{++}}\) is an algebraic basis of the space of central representative functions on \(\mathbf{G}\) and an orthonormal basis of the space \(\mathrm{ZL}^2(\mathrm{G})\) of classes of square-integrable central functions on \(\mathbf{G}\) .

Thus, from Prop. 3 and Th. 2 we deduce:

COROLLARY 3. The map which associates to any continuous central function \(\varphi\) on G the function \(w(\mathrm{G})^{-1/2}\mathrm{J}(\rho)(\varphi|\mathrm{T})\) induces an isomorphism from the space of central representative functions on G to the space of anti-invariant elements of \(\Theta(\mathrm{T})\) ; it extends by continuity to an isomorphism of Hilbert spaces from \(\mathrm{ZL}^2(\mathrm{G})\) to the subspace of anti-invariant elements of \(\mathrm{L}^2(\mathrm{T})\) .

COROLLARY 4. Let \(\varphi\) be a continuous central function on G. Then,

\[
\int_ {\mathrm{G}} \overline {{\chi_ {\lambda} (g)}} \varphi (g) d g = \int_ {\mathrm{T}} \overline {{\lambda (t)}} \prod_ {\alpha > 0} (1 - \alpha (t) ^ {- 1}) \varphi (t) d t = \int_ {\mathrm{T}} \overline {{\lambda \rho (t)}}. \varphi (t) \mathrm{J} (\rho) (t) d t.
\]

Indeed, by Lemma 4 and Th. 2,

\[
\begin{array}{c} \int_ {\mathrm{G}} \overline {{\chi_ {\lambda} (g)}} \varphi (g) d g = \frac {1}{w (\mathrm{G})} \int_ {\mathrm{T}} \overline {{\chi_ {\lambda} (t) \mathrm{J} (\rho) (t)}} (\varphi (t) \mathrm{J} (\rho) (t)) d t \\ = \frac {1}{w (\mathrm{G})} \int_ {\mathrm{T}} \overline {{\mathrm{J} (\lambda \rho) (t)}} \varphi (t) \mathrm{J} (\rho) (t) d t. \end{array}
\]

But the function \(t \mapsto \varphi(t)J(\rho)(t)\) is anti-invariant and \(\frac{1}{w(G)}J(\lambda\rho)\) is the orthogonal projection of \(\lambda\rho\) onto the subspace of anti-invariant elements of \(L^{2}(T)\) , so

\[
\frac {1}{w (\mathrm{G})} \int_ {\mathrm{T}} \overline {{\mathrm{J} (\lambda \rho) (t)}} \varphi (t) \mathrm{J} (\rho) (t) d t = \int_ {\mathrm{T}} \overline {{\lambda \rho (t)}} \varphi (t) \mathrm{J} (\rho) (t) d t;
\]

finally, by formula (3), we have \(\overline{\rho(t)}\mathrm{J}(\rho)(t) = \prod_{\alpha >0}(1 - \alpha (t)^{-1})\) , hence the corollary.

Remarks. 1) For all \(w \in \mathrm{W}\) , put \(\rho_w = {}^w\rho /\rho\) ; then

\[
\sum_ {w} \varepsilon (w) \rho_ {w} = \prod_ {\alpha > 0} (1 - \alpha^ {- 1}) = \rho^ {- 2} \prod_ {\alpha > 0} (\alpha - 1).\tag{6}
\]

If \(t\) is a regular element of \(\mathrm{T}\) , we deduce from (5) that

\[
\chi_ {\lambda} (t) = \frac {\sum_ {w} \varepsilon (w) ^ {w} \lambda (t) \rho_ {w} (t)}{\sum_ {w} \varepsilon (w) \rho_ {w} (t)} = \frac {\sum_ {w} \varepsilon (w) ^ {w} \lambda (t) \rho_ {w} (t)}{\prod_ {\alpha > 0} (1 - \alpha (t) ^ {- 1})}.\tag{7}
\]

Note that \(\rho_{w}\) is a linear combination of roots with integer coefficients, and hence belongs to X(T) even if we do not assume that \(\rho \in \mathrm{X}(T)\) . It follows that formula (7) is valid without the assumption that \(\rho \in \mathrm{X}(T)\) : indeed, to prove this we replace G by a suitable connected covering, and are then reduced to Cor. 2.

\begin{enumerate}
\def\labelenumi{\arabic{enumi})}
\setcounter{enumi}{1}
\item
  Similarly, the first equality of Cor. 4 remains valid without the assumption that \(\rho \in \mathrm{X}(\mathrm{T})\) .
\item
  We can deduce Th. 2 from the analogous infinitesimal statement (Chap. VIII, §9, no. 1, Th. 1); the same is true for Th. 3 of the next number (which is the analogue of Th. 2 of loc. cit., no. 2).
\end{enumerate}

\subsection{DEGREE OF IRREDUCIBLE REPRESENTATIONS}

We now return to the additive notation for the group \(X(T)\) and we no longer assume that \(\rho\) belongs to \(X(T)\) .

THEOREM 3. The dimension of the space of an irreducible representation of G of highest weight \(\lambda\) is given by

\[
\chi_ {\lambda} (e) = \prod_ {\alpha \in \mathrm{R} _ {+}} \frac {\langle \lambda + \rho , K _ {\alpha} \rangle}{\langle \rho , K _ {\alpha} \rangle}.
\]

Put \(\gamma = \frac{1}{2}\sum_{\alpha >0}K_{\alpha}\) , so \(\delta (\alpha)(\gamma) = 2\pi i\) for every simple root \(\alpha\) (Chap. VI, §1, no. 10, Prop. 29). The line \(\mathbf{R}\gamma\) is not contained in any of the hyperplanes \(\mathrm{Ker}\delta (\alpha)\) , so \(\exp (z\gamma)\) is a regular element of G for all sufficiently small \(z\in \mathbf{R}^*\) ; for all \(\mu \in \mathrm{X(T)}\) and all \(z\in \mathbf{R}\) , we have

\[
\mathrm{J} (\mu) (\exp (z \gamma)) = \sum_ {w \in \mathrm{W}} \varepsilon (w) e ^ {z \delta (\mu) (w ^ {- 1} \gamma)}.
\]

Let us accept provisionally the following lemma:

Lemma 5. We have

\[
\mathrm{J} (\mu) (\exp (z \gamma)) = e ^ {z \delta (\mu) (\gamma)} \prod_ {\alpha > 0} (1 - e ^ {- z \delta (\mu) (K _ {\alpha})}).
\]

Thus, the function \(\mathrm{J}(\mu)(\exp (z\gamma))\) is the product of a function that tends to 1 when \(z\) tends to 0 and of

\[
z ^ {\mathrm{N}} \prod_ {\alpha > 0} \delta (\mu) (K _ {\alpha}) = (2 \pi i z) ^ {\mathrm{N}} \prod_ {\alpha > 0} \langle \mu , K _ {\alpha} \rangle
\]

where N = Card R+.

Assume first of all that \(\rho \in \mathrm{X(T)}\) ; applying Cor. 2 of Th. 2 we see that, as \(z\) tends to 0, \(\chi_{\lambda}(z\gamma)\) tends to

\[
\prod_ {\alpha > 0} \langle \lambda + \rho , K _ {\alpha} \rangle / \prod_ {\alpha > 0} \langle \rho , K _ {\alpha} \rangle ,
\]

hence the theorem in this case.

In the general case, it suffices to remark that, in proving Th. 3, we can always replace G by a suitable connected covering and thus reduce to the preceding case.

We now prove Lemma 5. Let \(z \in \mathbf{C}\) ; denote by \(\varphi_z\) the map from \(\mathfrak{t}\) to the \(\mathbf{C}\) -algebra \(\operatorname{Map}(\mathrm{X}(\mathrm{T}), \mathbf{C})\) of maps from \(\mathrm{X}(\mathrm{T})\) to \(\mathbf{C}\) that associates to \(H \in \mathfrak{t}\) the map

\[
\varphi_ {z} (H): \mu \mapsto \mu (\exp z H) = e ^ {z \delta (\mu) (H)}.
\]

We have \(\varphi_z(H + H') = \varphi_z(H)\varphi_z(H')\) , so there exists a homomorphism of rings

\[
\psi_ {z}: \mathbf {Z} [ \mathfrak {t} ] \to \operatorname{Map} (\mathrm{X} (\mathrm{T}), \mathbf {C})
\]

such that \(\psi_z(e^H)(\mu) = e^{z\delta (\mu)(H)}\) . On the other hand, by Chap. VI, §3, no. 3, Prop. 2, we have the following relation in \(\mathbf{Z}[t]\) :

\[
\sum_ {w \in \mathrm{W}} \varepsilon (w) e ^ {w \gamma} = e ^ {\gamma} \prod_ {\alpha > 0} (1 - e ^ {- K _ {\alpha}}).
\]

Applying the homomorphism \(\psi_z\) , and using the equality

\[
\psi_ {z} (e ^ {w \gamma}) (\mu) = e ^ {z \delta (\mu) (w \gamma)} = e ^ {z \delta (w ^ {- 1} \mu) (\gamma)},
\]

we deduce the stated formula.

COROLLARY 1. Let \(\| \|\) be a norm on \(\mathrm{X(T)}\otimes \mathbf{R}\) . For all \(\lambda \in \mathrm{X}_{++}\) , let \(d(\lambda)\) be the dimension of the space of an irreducible representation of \(\mathrm{G}\) of highest weight \(\lambda\) .

\(a) \sup_{\lambda \in X_{++}} d(\lambda) / \| \lambda + \rho \|^{N} < \infty, \text{ where } N = 1/2 (\dim G - \dim T).\)

\begin{enumerate}
\def\labelenumi{\alph{enumi})}
\setcounter{enumi}{1}
\item
  If G is semi-simple, \(\inf_{\lambda\in X_{++}}d(\lambda)/\|\lambda+\rho\|>0\) .
\item
  For all \(\alpha \in \mathrm{R}_{+}\) , there exists \(\mathrm{A}_{\alpha} > 0\) with \(|\langle \lambda + \rho, K_{\alpha} \rangle| \leq \mathrm{A}_{\alpha} \| \lambda + \rho \|\) , hence \(d(\lambda) / \| \lambda + \rho \|^{\mathrm{N}} \leq \prod_{\alpha > 0} \mathrm{A}_{\alpha} / \langle \rho, K_{\alpha} \rangle\) .
\item
  Assume that G is semi-simple, denote by \(\beta_{1},\ldots,\beta_{r}\) the simple roots and put \(N_{i}=K_{\beta_{i}}\) . Then
\end{enumerate}

\[
d (\lambda) \geq \prod_ {i = 1} ^ {r} \frac {\langle \lambda + \rho , N _ {i} \rangle}{\langle \rho , N _ {i} \rangle} = \prod_ {i = 1} ^ {r} \langle \lambda + \rho , N _ {i} \rangle ;
\]

since \(\langle \lambda +\rho ,N_i\rangle \geq \langle \rho ,N_i\rangle = 1\) , this implies that

\[
d (\lambda) \geq \sup _ {i} | \langle \lambda + \rho , N _ {i} \rangle |.
\]

If G is semi-simple, \(x \mapsto \sup |\langle x, N_{i} \rangle|\) is a norm on \(\mathrm{X(T)} \otimes \mathbf{R}\) , necessarily equivalent to the given norm, hence b).

COROLLARY 2. Assume that G is semi-simple; let d be an integer. The set of classes of representations of G of dimension \(\leq d\) is finite.

Cor. 1 b) implies that the set \(X_{d}\) of elements \(\lambda\) of \(X_{++}\) such that \(d(\lambda) \leq d\) is finite. For all \(\lambda\) in \(X_{d}\) , let \(V_{\lambda}\) be an irreducible representation of highest weight \(\lambda\) ; every representation of dimension \(\leq d\) is isomorphic to a direct sum \(\bigoplus_{\lambda \in X_{d}} V_{\lambda}^{n_{\lambda}}\) , with \(n_{\lambda} \leq d\) , hence the corollary.

\subsection{CASIMIR ELEMENTS}

By Prop. 3 of §1, no. 3, there exist negative symmetric bilinear forms on g, separating and invariant under Ad(G) (if G is semi-simple, we can take the Killing form of g, for example). Let F be such a form. Recall (Chap. I, §3, no. 7) that the Casimir element associated to F is the element \(\Gamma\) of the centre of the enveloping algebra U(g) such that, for any basis \((e_{i})\) of g satisfying \(\mathrm{F}(e_{i}, e_{j}) = -\delta_{ij}\) , we have \(\Gamma = -\sum e_{i}^{2}\) .

In the remainder of this chapter we shall call the Casimir elements of G the elements of \(\mathrm{U}(\mathfrak{g})\) obtained in this way from negative separating invariant symmetric bilinear forms on g. If \(\Gamma\) is a Casimir element of G and if \(\tau: G \to \mathbf{GL}(V)\) is an irreducible representation of G, the endomorphism \(\Gamma_{V}\) of V is a homothety (Algebra, Chap. VIII, §3, no. 2, Th. 1), whose ratio we shall denote by \(\tilde{\Gamma}(\tau)\) .

PROPOSITION 4. Let \(\Gamma\) be the Casimir element of G.

\begin{enumerate}
\def\labelenumi{\alph{enumi})}
\item
  If \(\tau\) is an irreducible representation of \(\mathbf{G}\) , \(\tilde{\Gamma}(\tau)\) is real and positive. If \(\tau\) is not the trivial representation, \(\tilde{\Gamma}(\tau) > 0\) .
\item
  There exists a unique quadratic form \(\mathbf{Q}_{\Gamma}\) on \(\mathrm{X(T)}\otimes \mathbf{R}\) such that, for every irreducible representation \(\tau\) of \(\mathbf{G}\) ,
\end{enumerate}

\[
\tilde {\Gamma} (\tau) = \mathrm{Q} _ {\Gamma} (\lambda + \rho) - \mathrm{Q} _ {\Gamma} (\rho),
\]

where \(\lambda\) is the highest weight of \(\tau\) . The form \(Q_{\Gamma}\) is positive, separating and invariant under W.

Let \(\mathrm{F}\) be a separating negative symmetric bilinear form on \(\mathfrak{g}\) defining \(\Gamma\) . Let \(\tau : \mathrm{G} \to \mathbf{GL}(\mathrm{V})\) be an irreducible representation of \(\mathrm{G}\) ; let \(\langle , \rangle\) be a Hilbert scalar product on \(\mathrm{V}\) invariant under \(\mathrm{G}\) (§1, no. 1), and \((e_i)\) a basis of \(\mathfrak{g}\) such that \(\mathrm{F}(e_i, e_j) = -\delta_{ij}\) . Then, for every element \(v\) of \(\mathrm{V}\) not invariant under \(\mathrm{G}\) we have

\[
\begin{array}{c} \tilde {\Gamma} (\tau) \langle v, v \rangle = \langle v, \Gamma_ {\mathrm{V}} (v) \rangle = - \sum_ {i} \langle v, (e _ {i}) _ {\mathrm{V}} ^ {2} v \rangle = \sum_ {i} \langle v, (e _ {i}) _ {\mathrm{V}} ^ {*} (e _ {i}) _ {\mathrm{V}} v \rangle \\ = \sum_ {i} \langle (e _ {i}) _ {\mathrm{V}} v, (e _ {i}) _ {\mathrm{V}} v \rangle > 0, \end{array}
\]

hence \(a\) ).

Let B be the inverse form on \(t_{C}^{*}\) of the restriction to \(t_{C}\) of the bilinear form on \(g_{C}\) induced by F by extension of scalars. By the Cor. of Prop. 7 of Chap. VIII, §6, no. 4, we have \(\tilde{\Gamma}(\tau) = \mathrm{B}(\delta(\lambda), \delta(\lambda + 2\rho))\) . Extend \(\delta : \mathrm{X(T)} \to t_{\mathbf{C}}^{*}\) to an R-linear map from \(\mathrm{X(T)} \otimes \mathbf{R}\) to \(t_{C}^{*}\) and let \(Q_{\Gamma}\) be the quadratic form \(x \mapsto \mathrm{B}(\delta(x), \delta(x))\) on \(\mathrm{X(T)} \otimes \mathbf{R}\) ; it is separating and invariant under W, and

\[
\tilde {\Gamma} (\tau) = \mathrm{B} (\delta (\lambda + \rho), \delta (\lambda + \rho)) - \mathrm{B} (\delta (\rho), \delta (\rho)) = \mathrm{Q} _ {\Gamma} (\lambda + \rho) - \mathrm{Q} _ {\Gamma} (\rho).
\]

We show that the form \(Q_{\Gamma}\) is positive. Indeed, if \(x \in \mathrm{X}(\mathrm{T}) \otimes \mathbf{Q}\) , the element \(\delta(x)\) of \(t_{C}^{*}\) takes purely imaginary values on t, hence real values on it; we conclude by remarking that, for \(y \in it\) , we have \(\mathrm{F}(y, y) \geq 0\) .

It remains to prove the uniqueness assertion in b). Let Q be a quadratic form on \(X(T) \otimes \mathbf{R}\) satisfying the required condition, and let \(\Phi\) (resp. \(\Phi_{\Gamma}\) ) be the bilinear form associated to Q (resp. \(Q_{\Gamma}\) ). For \(\lambda, \mu \in X_{++}\) , we have

\[
\begin{array}{l} \Phi (\lambda , \mu) = (Q (\lambda + \mu + \rho) - Q (\rho)) - (Q (\lambda + \rho) - Q (\rho)) - (Q (\mu + \rho) - Q (\rho)) \\ \qquad = \Phi_ {\Gamma} (\lambda , \mu). \end{array}
\]

Since \(\mathrm{X(T)_{++}}\) generates the R-vector space \(\mathrm{X(T)}\otimes\mathbf{R}\) , we have \(\Phi=\Phi_{\Gamma}\) , so \(Q=Q_{\Gamma}\) .

Remark. Let \(x \in g\) . There exists a strictly positive real number A such that, for every irreducible representation \(\tau : G \to \mathbf{GL}(V)\) and every Hilbert structure on V invariant under G,

\[
\| \mathrm{L} (\tau) (x) \| ^ {2} \leq \mathrm{A}. \tilde {\Gamma} (\tau).
\]

Indeed, with the notations in the preceding proof, we can choose the basis \((e_{i})\) of g so that \(x = ae_{1}, a \in R\) . Then, for \(v \in V\) , we have

\[
\langle x _ {\mathrm{V}} v, x _ {\mathrm{V}} v \rangle = | a | ^ {2} \langle e _ {1} v, e _ {1} v \rangle \leq | a | ^ {2} \tilde {\Gamma} (\tau) \langle v, v \rangle .
\]

\section{FOURIER TRANSFORM}

We retain the notations and conventions of the preceding paragraph.